% concatenated from complete_3APfree.tex
\documentclass[11pt]{article}
\usepackage{enumerate}
\usepackage{amsmath}
\usepackage{amssymb,latexsym}
\usepackage{amsthm}
\usepackage{color}
\usepackage[normalem]{ulem}
%\usepackage{anysize}
\usepackage{graphicx}
%\usepackage{wrapfig}
\usepackage{mathabx}
\usepackage{hyperref}
%\usepackage[notcite, notref]{showkeys}
\usepackage{amscd}
%\usepackage{showlabels}
\usepackage{todonotes}
\DeclareMathOperator{\Tr}{Tr}
\DeclareMathOperator{\N}{N}
\usepackage{comment}
\usepackage{adjustbox}
\usepackage{tikz}
\usepackage{caption}
\usepackage{pgfplots}
\newcommand{\vgeq}{\mathrel{\rotatebox{90}{$\geq$}}}
% \newcommand{\vneq}{\mathrel{\rotatebox{90}{$\neq$}}}

\textwidth 18cm
\textheight 22 cm
\topmargin -1cm
\oddsidemargin -1cm
\evensidemargin -1cm
\parskip 1mm
\setlength{\parindent}{0pt}
%\usepackage{lineno,xcolor}
%% Running line numbers:
%\linenumbers
%\setlength\linenumbersep{15pt}
%\renewcommand\linenumberfont{\normalfont\tiny\sffamily\color{gray}}
\numberwithin{table}{section}

%\title{Complete 3-\mathrm{AP}-free sets and related structures in vector spaces}
%\title{Small, complete subsets of groups without a 3-term arithmetic progression}
\title{Complete $3$-term arithmetic progression free sets of small size in vector spaces and other abelian groups}
\author{Bence Csajb\'ok, \footnote{Dipartimento di Meccanica, Matematica e Management, Politecnico di Bari, Via Orabona 4, I-70125 Bari, Italy. Current address: Department of Computer Science, ELTE Eötvös Loránd University, Budapest, Hungary. E-mail: {\tt bence.csajbok@ttk.elte.hu}}\quad Zolt\'an L\'or\'ant Nagy \footnote{ELTE Linear Hypergraphs  Research Group, ELTE Eötvös Loránd University, Budapest, Hungary. The author is supported by the Hungarian Research Grant (NKFIH) No. PD  134953. and No. K. 124950 and the University Excellence Fund of Eötvös Loránd University E-mail: {\tt nagyzoli@cs.elte.hu}}}
\date{}



\newcommand{\nul}{{\underline 0}}
\newcommand{\ve}{{\underline v}}
\newcommand{\uE}{{\textup E}}
\newcommand{\uL}{{\textup L}}
\newcommand{\uS}{{\textup S}}
\newcommand{\Zp}{{\mathbb Z}_p}
\newcommand{\Z}{{\mathbb Z}}
\newcommand{\cR}{{\mathcal R}}
\newcommand{\cA}{{\mathcal A}}
\newcommand{\cB}{{\mathcal B}}
\newcommand{\cT}{{\mathcal T}}
\newcommand{\cC}{{\mathcal C}}
\newcommand{\cK}{{\mathcal K}}
\newcommand{\cN}{{\mathcal N}}
\newcommand{\cM}{{\mathcal M}}
\newcommand{\cG}{{\mathcal G}}
\newcommand{\ts}{{\theta_s}}
\newcommand{\cE}{{\mathcal E}}
\newcommand{\cF}{{\mathcal F}}
\newcommand{\cP}{{\mathcal P}}
\newcommand{\cH}{{\mathcal H}}
\newcommand{\cD}{{\mathcal D}}
\newcommand{\cU}{{\mathcal U}}
\newcommand{\cL}{{\mathcal L}}
\newcommand{\Lb}{{\mathbb L}}
\newcommand{\wa}{{\widehat{\alpha}}}
\newcommand{\wg}{{\widehat{\gamma}}}
\newcommand{\cJ}{{\mathcal J}}
\newcommand{\cQ}{{\mathcal Q}}
\newcommand{\cS}{{\mathcal S}}
\newcommand{\cX}{{\mathcal X}}
\newcommand{\F}{{\mathbb F}}
\newcommand{\Fq}{\F_q}
\newcommand{\Fqt}{\F_{q^t}}
\newcommand{\K}{{\mathbb K}}
\newcommand{\GF}{\hbox{{\rm GF}}}
\newcommand{\st}{{\cal S}_t}
\newcommand{\so}{{\cal S}}
\newcommand{\sk}{{\cal S}_2}
\newcommand{\ad}{E}
\newcommand{\lmb}{\lambda}
\newcommand{\la}{\langle}
\newcommand{\ra}{\rangle}
%\renewcommand{\proof}{\noindent{\bf Proof.}\ }
%\newcommand{\Qed}{\hfill $\Box$ \medskip}
\renewcommand{\mod}{\hbox{{\rm mod}\,}}
\newcommand{\ZG}{\mathcal{Z}
	(\mathrm{\Gamma L})}
\newcommand{\G}{\mathrm{\Gamma L}}
\newcommand{\E}{\mathbb{E}}
\newcommand{\pP}{\mathbb{P}}
\newcommand{\sat}{\mathrm{sat}}
\newcommand{\ep}{\epsilon}
\newcommand{\vv}{\mathbf v}
\newcommand{\hs}{{\hat{\sigma}}}

\newtheorem{theorem}{Theorem}[section]
\newtheorem{problem}[theorem]{Problem}
\newtheorem{lemma}[theorem]{Lemma}
\newtheorem{corollary}[theorem]{Corollary}
\newtheorem{definition}[theorem]{Definition}
\newtheorem{proposition}[theorem]{Proposition}
\newtheorem{result}[theorem]{Result}
\newtheorem{example}[theorem]{Example}
\newtheorem{conjecture}[theorem]{Conjecture}
\newtheorem{observ}[theorem]{Observation}
\newtheorem{remark}[theorem]{Remark}
%\newtheorem{theorem}{\bf Theorem}[section]
\newtheorem{constr}[theorem]{\bf Construction}
\newtheorem{cor}[theorem]{\bf Corollary}
%\newtheorem{proposition}[theorem]{\bf Proposition}

\newtheorem{prop}[theorem]{\bf Proposition}
\newtheorem{conj}[theorem]{\bf Conjecture}


\newtheorem{notation}[theorem]{\bf Notation}


\theoremstyle{definition} 
\newtheorem{defi}[theorem]{\bf Definition}


\newcommand{\red}[2]{\textcolor[rgb]{1,0.1,0.1}{#1}}
\newcommand{\green}[2]{\textcolor[rgb]{0.2,0.9,0.1}{#1}}

\DeclareMathOperator{\tr}{Tr}
\DeclareMathOperator{\PG}{{PG}}
\DeclareMathOperator{\AG}{{AG}}
\DeclareMathOperator{\GL}{{GL}}
\DeclareMathOperator{\PGL}{{PGL}}
\DeclareMathOperator{\PGaL}{P\Gamma L}
\DeclareMathOperator{\Gal}{Gal}
\DeclareMathOperator{\GaL}{\Gamma L}


\newcommand{\corr}[2]{\textcolor[rgb]{1,0.1,0.1}{#1}}

\begin{document}
	\maketitle
	
	\begin{abstract}
		
		% Let $G$ denote an abelian group, written additively. A $3$-term arithmetic progression of $G$, $3$-$\mathrm{AP}$ for short, is a set of three distinct elements of $G$ of the form $g$, $g+d$, $g+2d$, for some $g,d \in G$. 
		
		A subset $S$ of an abelian group $G$ is called $3$-$\mathrm{AP}$ free if it does not contain a three term arithmetic progression. Moreover, $S$ is called complete $3$-$\mathrm{AP}$ free, if it is maximal w.r.t. set inclusion.
		% $S$ is called complete $3$-$\mathrm{AP}$ free if it is  $3$-$\mathrm{AP}$ free and the addition of any further element would violate the  $3$-$\mathrm{AP}$ free property. 
		%  In case $S$ is not necessarily $3$-$\mathrm{AP}$ free but for 
		% each $x \in
		% G \setminus S$ there is a $3$-$\mathrm{AP}$ of $G$ consisting of $x$ and two elements of $S$, we call $S$ saturating w.r.t. $3$-$\mathrm{AP}$s.
		%Moreover, $S$ is called complete $3$-$\mathrm{AP}$ free if $S$ is $3$-$\mathrm{AP}$ free and it is not contained in a larger $3$-$\mathrm{AP}$ free set. A set $S$ is $3$-$\mathrm{AP}$ saturating if for each $x \in G \setminus S$ there is a $3$-$\mathrm{AP}$ of $G$ consisting of $x$ and two elements of $S$. 
		%Clearly, $S$ is complete $3$-$\mathrm{AP}$ free if and only if it is $3$-$\mathrm{AP}$ free and $3$-$\mathrm{AP}$ saturating.
		One of the most central problems in additive combinatorics is to determine the maximal size of a $3$-$\mathrm{AP}$ free set, which is necessarily complete. 
		% A classical problem in finite geometry and coding theory is to find the minimum size of a complete cap in affine spaces, where a cap is a point set without a collinear triplet and in $\F_3^n$ caps and $3$-$\mathrm{AP}$ free sets are the same objects.
		% If $G=\mathbb{F}_3^n$ then the $3$-$\mathrm{AP}$ free sets are exactly the affine point sets without a collinear triplet. Point sets of $V=\mathbb{F}_q^n$ without a collinear triplet are called (affine) caps. 
		% A classical problem in finite geometry is to find small complete caps (a cap is complete if it is not contained in a larger cap) and small saturating sets (a point set $S$ is saturating if for each point $x \in V \setminus S$ there is a collinear triplet of $V$ consisting of $x$ and two elements of $S$).
		In this paper we are interested in the minimum size of complete $3$-$\mathrm{AP}$ free sets. We define and study saturation w.r.t. $3$-$\mathrm{AP}$s and present constructions of small complete $3$-$\mathrm{AP}$ free sets 
		and $3$-$\mathrm{AP}$ saturating sets 
		for several families of vector spaces and cyclic groups.
		%, whose size reaches the trivial lower bound up to a small multiplicative constant factor. 
		
		%present constructions of complete $3$-$\mathrm{AP}$ free sets and $3$-$\mathrm{AP}$ saturating sets whose size is close to the trivial lower bound. We are mostly interested in the case when $G$ is a finite vector space or a cyclic group.
	\end{abstract}
	
	%\bigskip
	%{\it AMS subject classification:}
	
	%\bigskip
	%{\it Keywords:}
	
	\section{Introduction}
	
	
	Studying the maximum possible size of a subset of a vector space  over a finite field which contain either no (non-trivial) solution to a given linear equation or not too many collinear points is a classical yet vibrant research area \cite{Alon, Bruyn,Eber, EG, KovacsNagy, FGR, Mimura, Lisa, Ilja}.
	The most notable examples are the Sidon sets and the so-called cap-set problem. The latter one is to determine the largest subset of $\F_3^n$ containing no complete line, or in other terms, no arithmetic progression of length 3 ($3$-$\mathrm{AP}$), or no collinear triplets. In general,  point sets of finite affine or projective spaces with no three in line are called caps. Recently Ellenberg and Gijswijt proved a breakthrough result regarding caps in $\F_3^n$ \cite{EG} building on the ideas of Croot, Lev and Pach \cite{CLP}, and they also proved that the size of a 
	$3$-$\mathrm{AP}$ free set of $\F_p^n$ is always bounded from above by $(p-\delta_p)^n$ for some small constant $\delta_p>0$ depending only on $p$, see also \cite{EPPP}.%\todo[inline]{itt elmondani mi a cap?}
	
	For the general case of abelian groups,  the maximum size of $3$-$\mathrm{AP}$ free sets  was discussed in the classical paper of Frankl, Graham and Rödl \cite{FGR}.   A nice and general overview on additive combinatorial and extremal results concerning arithmetic progressions is by Shkredov \cite{Ilja}.
	
	A finite point set (with respect to a property)  is called  \textit{complete} if it does not contained as a subset in a larger point set, satisfying the same property.
	In extremal problems described above, usually constructions of maximum size are in the center of attention. These are complete by definition.  However, in several cases, the whole spectra of sizes matters for complete structures, and the structure of smallest size in particular. For example, in the case of complete caps over $\PG(n,q)$, the point set is corresponding to the parity check matrix of a $q$-ary linear code with codimension $n+1$,
	Hamming distance $4$, and covering radius $2$, see \cite{Hirs}. 
	
	In this paper we investigate the less studied lower end of the spectrum of possible sizes of complete $3$-$\mathrm{AP}$ free sets, the minimum size. We discuss the minimum size in arbitrary abelian groups of odd order and highlight the case of finite vector spaces.
	
	
	%the size of small complete $3$-$\mathrm{AP}$ free subsets of vector spaces over finite fields, and some related problems.
	
	
	%This problem also has its analogues for  subsets of  the integers of an interval $[1,n]$ as well, see \cite{}.
	
	% \noindent Let $G$  denote %ilyenkor nem kell az s (ezt a kommentet most lehet hogy magamnak irom :))
	% an Abelian group, where the operation is the addition. Take some $g,d \in G$, then $\{g,g+d,g+2d\}$ is a set of three distinct elements of $G$ if and only if the order of $d$ is at least $3$.
	
	
	
	A $3$-term arithmetic progression of the abelian group $G$, $3$-$\mathrm{AP}$ for short,  is a set of three distinct elements of $G$ of the form $g,g+d,g+2d$, where $g,d\in G$. We will call $d$ the difference. If $d$ is the difference of a $3$-$\mathrm{AP}$ in $G$, then the order of $d$ is at least $3$. Let $\F_q$ denote the Galois field of $q$ elements, and $o_q(a)$ denotes the multiplicative order of $a \in\F_q$. In order to have a $3$-$\mathrm{AP}$ in $\F_q^n$ we need $q$ to be odd, so we will only consider this case. $A+B$ denotes the sumset $\{a+b :  a\in A, b\in B\}$ of sets $A$ and $B$,  while $A\dot+B$ denotes the restricted sumset where the summands must be distinct.
	
	
	\begin{definition}
		$A\subseteq G$ is called \textbf{$3$-$\mathrm{AP}$-free} if it does not contain a $3$-$\mathrm{AP}$ of $G$.
		Moreover, $A$ is \textbf{complete $3$-$\mathrm{AP}$-free} if it is $3$-$\mathrm{AP}$ free and  not contained in a larger $3$-$\mathrm{AP}$ free set. 
	\end{definition}
	\vspace{-0.2cm}
	The completeness of a $3$-$\mathrm{AP}$ free set can be interpreted via saturation as well: $S$ is complete $3$-$\mathrm{AP}$ free if $S$ is $3$-$\mathrm{AP}$ free and a saturating set w.r.t. $3$-$\mathrm{AP}$s.
	
	
	\begin{definition}
		For a subset $S\subseteq G$ we say that $S$ is\textbf{ $3$-$\mathrm{AP}$ saturating } or in other words, $S$ \textbf{$3$-$\mathrm{AP}$ saturates } $G$  if for each $x \in G \setminus S$  there is a $3$-$\mathrm{AP}$ of $G$ consisting of $x$ and two  elements of $S$. %Moreover, if for every $x \in G \setminus A$  there is a $3$-$\mathrm{AP}$ $\{x,y,z\}$ such that $x=2y-z$  for a pair $\{y,z\}\subset A$, then $A$ is called \textbf{ $3$-$\mathrm{AP}$* saturating.} 
		In a broader context, we say that $S$ \textbf{$3$-$\mathrm{AP}$ saturates } 
		a set $H\subset G$ if similar condition holds for the elements of $H\setminus S$.
	\end{definition}
	\vspace{-0.2cm}
	In this paper we will mostly consider the problem of $3$-$\mathrm{AP}$ saturation in groups of odd order. 
	%In this this case the following holds.
	\vspace{-0.2cm}
	\begin{definition}
		For $g\in G$ if $\ell\in \mathbb{Z}$ is positive, then $\ell g= \underbrace{g+g+\ldots+g}_{\ell} \in G$ and $(-\ell)g=~-~(\ell g)\in ~G$. 
		If the order of $G$ is odd then we define $\frac12 g$ as the unique element $x\in G$ such that $x+x=g$, that is,  $x=(k+1)g$, where the order of $g$ is $2k+1$.
	\end{definition}
	
	%\noindent Observe that { $3$-$\mathrm{AP}$* saturating} sets are also { $3$-$\mathrm{AP}$ saturating,} and $A$ is {complete $3$-$\mathrm{AP}$-free}  if $A$ is $3$-$\mathrm{AP}$ free and it saturates the $3$-$\mathrm{AP}$s of $G$.
	
	\begin{observ}\label{obs1}
		A set $A\subseteq G$ is $3$-$\mathrm{AP}$ saturating, iff  for every $x \in G \setminus A$  there is a $3$-$\mathrm{AP}$  consisting of $x$ and $\{a_1, a_2\}\subseteq A$ such that $(i)$ either  $x=2a_1-a_2$, or $(ii)$ $2x=a_1+a_2$.
		
		If $G$ has odd order then $(ii)$ is equivalent to   %\vspace{-0.2cm}
		%\begin{enumerate}[\rm(3)]
		%\item 
		$x=\frac12 a_1 + \frac12 a_2$.
		% \end{enumerate}
	\end{observ}
	\vspace{-0.3cm}
	Saturation and completeness with respect to  $3$-$\mathrm{AP}$s  were considered in integer sequences as well, see \cite{SandorCs, Fang}.\\
	The postage stamp problem seeks the greatest integer  $r=r_k$ such that  there exists a set $A_k$ of $k$ positive integers together with $0$ such that \vspace{-0.3cm} $$i\in A_k+A_k \ \ \mbox{ for all \  } i=0, 1, \ldots, r.$$ 
	Mrose \cite{Mrose} and independently,  Fried \cite{Fried}  showed that the  $r_k$ is at least $\frac{2}{7}k^2+O(k)$. %\todo{ezt lehetne egzaktul}
	A set of non-negative integers $A$ is called a basis of order two if $A+A=\mathbb{N}$ holds for the sumset. 
	Note that a variation of Mrose's construction can be extended to an additive $2$-basis of $\mathbb{N}$, see \cite{Habsi, Hof}. Such constructions lead to 
	small sets $S$ of $\F_p$ for which every element of $\F_p$ is an arithmetic mean of a pair of elements from $S$, hence $S$ saturates the $3$-$\mathrm{AP}$s. We will discuss this in Section 4.
	
	These constructions provide further motivations to introduce the saturation with respect to a set of (coefficient) vectors. %Indeed, the above sequences are formed to avoid the solutions of a single linear equation $x=2a-a'$.
	
	% \begin{definition}[$W$-avoiding and $W$-saturating sets]
	%     Let $W$ denote a set of vectors from $\F_q^*\times \F_q^*$.
	%     We define $W$-avoiding and $W$-saturating sets in $\F_q^n$ as follows. 
	
	%     $A\subseteq \F_q^n$ is  \textbf{$W$-avoiding}, if there is no $w=(\lambda_1, \lambda_2)\in W$ such that $a=\lambda_1a'+\lambda_2a''$ has a nontrivial solution 
	%     with $a,a',a''$ pairwise distinct vectors of $A$. 
	%     %$(a, a', a'')\in A^3.$ \textcolor{blue}{A solution is trivial if $a$ coincides with $a'$ or $a''$.}
	%     %\textcolor{red}{Ezt most még nem latom miért ezek a trivialisak és miért nem mondjuk az, hogy $\{a,a',a''\}$ mérete kisebb mint $3$.} 
	
	
	% $A\subseteq \F_q^n$ is  \textbf{$W$-saturating} in $\F_q^n$ if for each    $x\in \F_q^n\setminus A$
	% there exists a $w=(\lambda_1, \lambda_2)\in W$ such that $x=\lambda_1a'+\lambda_2a''$ for a pair $(a', a'')\in A^2$, $a'\neq a''$. 
	% \end{definition}
	
	For a subset $S$ of elements of a group  $G$ (written additively) we will write $S^*$ to denote $S$ minus the neutral element. 
	
	\begin{definition}[$W$-avoiding and $W$-saturating sets]
		\label{defavoid}
		Let $W$ denote a set of vectors from $\F_q^*\times \F_q^*$. 
		We define $W$-avoiding and $W$-saturating sets in $\F_q^n$ as follows. 
		
		$A\subseteq \F_q^n$ is  \textbf{$W$-avoiding}, if there is no $w=(\lambda_1, \lambda_2)\in W$ such that $a=\lambda_1a'+\lambda_2a''$ has a non-trivial solution 
		with $a,a',a''$ pairwise distinct vectors of $A$. 
		%$(a, a', a'')\in A^3.$ \textcolor{blue}{A solution is trivial if $a$ coincides with $a'$ or $a''$.}
		%\textcolor{red}{Ezt most még nem latom miért ezek a trivialisak és miért nem mondjuk az, hogy $\{a,a',a''\}$ mérete kisebb mint $3$.} 
		
		
		$A\subseteq \F_q^n$ is  \textbf{$W$-saturating} in $\F_q^n$ if for each    $x\in \F_q^n\setminus A$
		there exists a $w=(\lambda_1, \lambda_2)\in W$ such that $x=\lambda_1a'+\lambda_2a''$ for a pair $(a', a'')\in A^2$, $a'\neq a''$. 
		
		$A \subseteq \F_q^n$ is \textbf{complete $W$-avoiding} if it is $W$-avoiding and $W$-saturating. 
		
	\end{definition}
	
	
	
	If $W$ consists of a single vector $W=\{w\}$, we omit the brackets for brevity.
	
	
	\begin{definition}[Avoiding and saturating sets in groups]
		
		Let $W$ denote a subset of $\mathbb{Z}\times \mathbb{Z}$.   
		We define $W$-avoiding, $W$-saturating and complete $W$-avoiding sets in  abelian groups $G$ similarly to Definition \ref{defavoid}. We will be mostly interested in the cases when $W$ is an element of $\{(2,-1), (1,1),(1,-1)\}$. 
		
		If $G$ has odd order then we also define $(\frac12,\frac12)$-saturating sets.
		
	\end{definition}
	
	\begin{remark}
		\label{r1}
		For a subset $A\subseteq \F_q^n$ and for $(\lambda_1,\lambda_2)\in \F_q^*\times \F_q^*$ the following properties are equivalent: $(i)$ $A$ is $(\lambda_1,\lambda_2)$-avoiding, $(ii)$ $A$ is $(1/\lambda_1,-\lambda_2/\lambda_1)$-avoiding, $(iii)$ $A$ is $(1/\lambda_2,-\lambda_1/\lambda_2)$-avoiding.
		
		In a group $G$ the following properties are equivalent: $(i)$ $A$ is $3$-$\mathrm{AP}$ free $(ii)$ $A$ is $(2,-1)$-avoiding, $(iii)$ there is no three pairwise distinct elements $x,y,z\in A$ such that $2x=y+z$. (If the order of $G$ is odd, then it is equivalent to say that $A$ is $(\frac12, \frac12)$-avoiding.)
		
		Similarly, for $A \subseteq G \setminus \{0\}$ the following properties are equivalent in any abelian group $G$: $(i)$ $A$ is restricted sum-free, i.e., $(A\dot+A)\cap A= \emptyset$, $(ii)$ $A$ is $(1,1)$-avoiding, $(iii)$ $A$ is $(1,-1)$-avoiding.
	\end{remark}
	
	% \begin{example}
	%  $3$-$\mathrm{AP}$ free property is equivalent to $W$-avoiding with   $W=\{(2,-1), (1/2, 1/2)\}$.  
	% \end{example}
	
	\vspace{-0.3cm}
	Note that Definition \ref{defavoid} can be extended naturally to any set of vectors of $\bigcup_{t=2}^\infty (\F_q^*)^t$.
	
	Our main (but not only) focus will be the case $W=\{(2,-1)\}$
	for its correspondence to  $3$-$\mathrm{AP}$s, see Observation \ref{obs1}, and in general, the case when $W$ consists of a single vector.
	
	
	We will also apply the fact that the field $\F_{q^n}$ is itself a vector space over $\F_{q^h}$ for  $h \mid n$. This will enable us to alter the dimension of the underlying structure at times, which will provide improvements on the estimates. If $q=p^r$, $p$ prime, and $W \subseteq \F_p^* \times \F_p^*$, then studying $W$-saturating and $W$-avoiding sets in the vector space $\F_q^n$ is equivalent to study the same questions in the elementary abelian group $\F_p^{rn}$.
	
	Now we introduce the functions we wish to study.
	
	\begin{definition}
		For a group $G$ we define the following: \vspace{-0.2cm}
		\begin{enumerate}[\rm(1)]\vspace{-0.2cm}
			\item  Let $a(3-\mathrm{AP}, G)$ denote the minimum size of a complete $3$-$\mathrm{AP}$ free set of $G$.\vspace{-0.2cm}
			\item  Let $\sat(3-\mathrm{AP}, G)$ denote the minimum size of a $3$-$\mathrm{AP}$ saturating set of $G$.\vspace{-0.2cm}
			\item  Let $a(W, G)$ denote the minimum size of a complete $W$-avoiding set of $G$. If there is no $W$-avoiding set of $G$, then put $a(W,G)=\infty$.\vspace{-0.2cm}
			\item  Let $\sat(W, G)$ denote the minimum size of a $W$-saturating set of $G$.\vspace{-0.2cm}
		\end{enumerate}
		% When $G=\F_q^n$ then we will simply write $a(3-\mathrm{AP}, q^n)$, $\sat(3-\mathrm{AP}, q^n)$, $a(W, q^n)$ and $\sat(W, q^n)$, respectively.
	\end{definition}
	
	Observe that in $\mathbb{Z}_5$ there are no complete $(2,-1)$-avoiding sets.
	
	\begin{example} As an example, we show complete $3$-AP sets for $G=\F_3^2$ and $G=\F_5^2$ in Figure \ref{fig:kicsipelda}. These are of minimum size, c.f. Lemma \ref{satbound}. We note in advance that in $\F_q^2$ we can always find complete $3$-AP sets of size $q$ when $-2$ is not a square element in $\F_q$, c.f. Theorem \ref{main1}.
	\end{example}
	
	% \definecolor{cqcqcq}{rgb}{0.7529411764705882,0.7529411764705882,0.7529411764705882}
	% \definecolor{ffffff}{rgb}{1.,1.,1.}
	% \begin{tikzpicture}[line cap=round,line join=round,x=1.0cm,y=1.0cm]
	% \begin{axis}[x=1.0cm,y=1.0cm, axis lines=middle,ymajorgrids=true,xmajorgrids=true,
	% xmin=-3.3029237069943105,
	% xmax=11.389629833048145,
	% ymin=-6.028382926531811,
	% ymax=6.571274497734802,
	% xtick={-2.0,0.0,...,10.0},
	% ytick={-6.0,-4.0,...,6.0}]
	% \clip(-3.3029237069943105,-6.028382926531811) rectangle (11.389629833048145,6.571274497734802);
	% \draw [line width=2.pt,color=cqcqcq] (0.,0.)-- (0.,2.);
	% \draw [line width=2.pt,color=cqcqcq] (0.,0.)-- (2.,0.);
	% \draw [line width=2.pt,color=cqcqcq] (2.,0.)-- (2.,2.);
	% \draw [line width=2.pt,color=cqcqcq] (0.,2.)-- (2.,2.);
	% \draw [line width=2.pt,color=cqcqcq] (0.,1.)-- (2.,1.);
	% \draw [line width=2.pt,color=cqcqcq] (1.,0.)-- (1.,2.);
	% \draw (0.2697979451886749,-0.06680005129156803) node[anchor=north west] {$\mathbb{F}_3\times \mathbb{F}_3$};
	% \draw [line width=2.pt,color=cqcqcq] (4.,0.)-- (4.,4.);
	% \draw [line width=2.pt,color=cqcqcq] (4.,4.)-- (8.,4.);
	% \draw [line width=2.pt,color=cqcqcq] (8.,4.)-- (8.,0.);
	% \draw [line width=2.pt,color=cqcqcq] (8.,0.)-- (4.,0.);
	% \draw [line width=2.pt,color=cqcqcq] (4.,1.)-- (8.,1.);
	% \draw [line width=2.pt,color=cqcqcq] (4.,2.)-- (8.,2.);
	% \draw [line width=2.pt,color=cqcqcq] (4.,3.)-- (8.,3.);
	% \draw [line width=2.pt,color=cqcqcq] (5.,4.)-- (5.,0.);
	% \draw [line width=2.pt,color=cqcqcq] (6.,4.)-- (6.,0.);
	% \draw [line width=2.pt,color=cqcqcq] (7.,4.)-- (7.,0.);
	% \draw (5.3012047689729975,-0.06680005129156803) node[anchor=north west] {$\mathbb{F}_5\times \mathbb{F}_5$};
	% \begin{scriptsize}
	% \draw [fill=black] (1.,0.) circle (2.5pt);
	% \draw [fill=ffffff] (2.,0.) circle (2.5pt);
	% \draw [fill=black] (0.,1.) circle (2.5pt);
	% \draw [fill=ffffff] (0.,2.) circle (2.5pt);
	% \draw [fill=black] (1.,1.) circle (2.5pt);
	% \draw [fill=ffffff] (1.,2.) circle (2.5pt);
	% \draw [fill=ffffff] (2.,1.) circle (2.5pt);
	% \draw [fill=ffffff] (2.,2.) circle (2.5pt);
	% \draw [fill=black] (0.,0.) circle (2.5pt);
	% \draw [fill=black] (4.,0.) circle (2.5pt);
	% \draw [fill=ffffff] (5.,0.) circle (2.5pt);
	% \draw [fill=ffffff] (6.,0.) circle (2.5pt);
	% \draw [fill=ffffff] (7.,0.) circle (2.5pt);
	% \draw [fill=ffffff] (8.,0.) circle (2.5pt);
	% \draw [fill=ffffff] (4.,1.) circle (2.5pt);
	% \draw [fill=black] (5.,1.) circle (2.5pt);
	% \draw [fill=ffffff] (6.,1.) circle (2.5pt);
	% \draw [fill=ffffff] (7.,1.) circle (2.5pt);
	% \draw [fill=black] (8.,1.) circle (2.5pt);
	% \draw [fill=ffffff] (4.,2.) circle (2.5pt);
	% \draw [fill=ffffff] (5.,2.) circle (2.5pt);
	% \draw [fill=ffffff] (6.,2.) circle (2.5pt);
	% \draw [fill=ffffff] (7.,2.) circle (2.5pt);
	% \draw [fill=ffffff] (8.,2.) circle (2.5pt);
	% \draw [fill=ffffff] (4.,3.) circle (2.5pt);
	% \draw [fill=ffffff] (4.,4.) circle (2.5pt);
	% \draw [fill=ffffff] (5.,3.) circle (2.5pt);
	% \draw [fill=ffffff] (6.,3.) circle (2.5pt);
	% \draw [fill=ffffff] (7.,3.) circle (2.5pt);
	% \draw [fill=ffffff] (8.,3.) circle (2.5pt);
	% \draw [fill=ffffff] (5.,4.) circle (2.5pt);
	% \draw [fill=black] (6.,4.) circle (2.5pt);
	% \draw [fill=black] (7.,4.) circle (2.5pt);
	% \draw [fill=ffffff] (8.,4.) circle (2.5pt);
	% \end{scriptsize}
	% \end{axis}
	% %\caption{Complete $3$-AP free sets of minimum size in $\F_3^2$ and in $\F_5^2$.}
	% \end{tikzpicture}
	% %\begin{figure}
	%  %   \centering
	%   %  \includegraphics[width=0.5\linewidth]{}
	%   %  \caption{Caption}
	%    % \label{fig:kicsipelda}
	% %\end{figure}
	
	\begin{figure}[htbp] % htbp = here, top, bottom, page (placement preference)
		\centering % Center the figure content
		
		% Define styles for points (optional, but good practice)
		\tikzstyle{small dot}=[circle, fill, inner sep=1.5pt]
		\tikzstyle{large dot}=[circle, fill, inner sep=3pt] % Use inner sep for size, or radius
		
		% --- First Grid (3x3) ---
		\begin{minipage}{0.45\textwidth} % Adjust width as needed
			\centering
			\begin{tikzpicture}
				% Define grid boundaries (coordinates from 0 to 2)
				\def\maxCoord{2}
				
				% Draw the grid lines
				\draw[step=1, gray, very thin] (0,0) grid (\maxCoord,\maxCoord);
				
				% Draw all grid points with small bullets
				\foreach \x in {0,...,\maxCoord} {
					\foreach \y in {0,...,\maxCoord} {
						\node[small dot] at (\x,\y) {};
					}
				}
				
				% Enlarge specific points
				\node[large dot] at (0,0) {};
				\node[large dot] at (0,1) {};
				\node[large dot] at (1,1) {};
				\node[large dot] at (1,0) {};
				
			\end{tikzpicture}
			%\caption*{3x3 Grid} % Optional caption for the subfigure
		\end{minipage}
		\hfill % Add horizontal space between the minipages
		% --- Second Grid (5x5) ---
		\begin{minipage}{0.45\textwidth} % Adjust width as needed
			\centering
			\begin{tikzpicture}
				% Define grid boundaries (coordinates from 0 to 4)
				\def\maxCoord{4}
				
				% Draw the grid lines
				\draw[step=1, gray, very thin] (0,0) grid (\maxCoord,\maxCoord);
				
				% Draw all grid points with small bullets
				\foreach \x in {0,...,\maxCoord} {
					\foreach \y in {0,...,\maxCoord} {
						\node[small dot] at (\x,\y) {};
					}
				}
				
				% Enlarge specific points
				\node[large dot] at (0,0) {};
				\node[large dot] at (1,1) {};
				\node[large dot] at (2,4) {};
				\node[large dot] at (3,4) {};
				\node[large dot] at (4,1) {};
				
			\end{tikzpicture}
			%\caption*{5x5 Grid} % Optional caption for the subfigure
		\end{minipage}
		
		\caption{Complete $3$-AP free sets of minimum size in $\F_3^2$ and in $\F_5^2$.}
		\label{fig:kicsipelda}
		% Label for referencing the figure
	\end{figure}
	
	\begin{remark}\label{rem}
		Any  complete $3$-$\mathrm{AP}$ free set is $3$-$\mathrm{AP}$ saturating, any complete $W$-avoiding set is $W$-saturating by definition, thus 
		\[a(W, G)\geq \sat(W, G).\]  %\begin{equation}
		%     a(3-\mathrm{AP}, q^n)\geq \sat(3-\mathrm{AP}, q^n) \quad \mbox{and} \quad a(W, q^n)\geq \sat(W, q^n) 
		%\end{equation} 
		
		Since %$3$-$\mathrm{AP}$ free sets are also  $(2,-1)$-avoiding sets, and
		$(2,-1)$-saturating sets are clearly  satisfying the $3$-$\mathrm{AP}$ saturating property in view of Observation \ref{obs1}, we have
		\begin{equation}\label{eq:triv}
			%a( (2,-1), q^n) \geq  a(3-\mathrm{AP}, q^n)  \mbox{\ and  \ }
			\begin{array}{ccc}
				\sat( (2,-1), G) & \geq & \sat(3-\mathrm{AP}, G)\\
				& & \\
				\vgeq &  & \vgeq\\
				& & \\
				a((2,-1),G) & \geq & a(3-\mathrm{AP}, G)\\
			\end{array}
		\end{equation}
	\end{remark}
	
	\vspace{-0.3cm}
	Our main results are as follows.
	
	
	
	\begin{theorem}\label{main3AP}
		Let $p$ be an odd prime and $k$ a positive integer. Then we have
		\begin{enumerate}[\rm(1)] \vspace{-0.2cm}
			\item 
			\[\sqrt{2/3}\cdot p^{2k-1}<a(3-\mathrm{AP}, \F_p^{4k-2})\leq p^{2k-1},\]  provided that  $-2$ is not a square element  in $\F_p$.  Also, \vspace{-0.2cm}
			\item %\[\sqrt{2/3}\cdot p^{2k+r/2-1}<a(3-\mathrm{AP}, \F_p^{4k+r-2})\leq c_p^r\cdot p^{2k+r/2-1},\]
			%provided that $-2$ is not a square element  in $\F_p$ and $\frac23(4^n-1) < p<4^n$ holds for some positive integer $n$. Here $c_p<...$ is a constant depending only on $p$ and $r\in\{1,2,3\}$.
			\[\sqrt{2/3}\cdot p^{n/2}<a(3-\mathrm{AP}, \F_p^{n})\leq a((2,-1),\F_p)^n.\]
		\end{enumerate}
	\end{theorem}
	\vspace{-0.3cm}
	Observe that $(2)$ of Theorem \ref{main3AP} motivates the study of 
	$a((2,-1),\F_p)$, or in general $a((2,-1),\Z_m)$, where $\Z_m$ is the cyclic group of order $m$.
	
	
	% nincs minden p-re még $a((2,-1),\F_p)$, ezen gondolkozni.
	
	
	\begin{theorem}\label{maingroups} Let $m$ denote a positive odd integer. Then
		\begin{enumerate}[\rm(1)] \vspace{-0.2cm}
			\item     $\sat((2,-1), \mathbb{Z}_m)<\sqrt{c_m\cdot m}$ where $c_m\in [1,3]$ is a constant depending only on $m$.
			\vspace{-0.2cm}
			\item %\textcolor{blue}{   
			$a((2,-1), \mathbb{Z}_m)<\sqrt{c_m \cdot m}$ for $c_m\in [1, 1.5]$, a constant depending only on $m$, provided that $\frac23(4^n-1) < m<4^n$ holds for some positive integer $n$.\vspace{-0.2cm}
			\item $\sqrt{2m}-0.5<\sat((1/2, 1/2), \mathbb{Z}_m)\leq (\sqrt{3.5}+o(1))\sqrt{m}\approx 1.87\sqrt{m}.$
			\vspace{-0.2cm}
			\item  $a((2,-1), \mathbb{Z}_m)=\lceil\sqrt{m}\rceil$, provided that $m$ is if form $m=2^{2t}+2^t+1$ for some positive integer $t$.%\textcolor{blue}{Singer }
		\end{enumerate}
	\end{theorem}
	\vspace{-0.3cm}
	The close connection between the $\sat$ function and the size of the complete $3$-AP-free sets, see Remark \ref{rem}, motivates the theorem below.
	
	\begin{theorem}\label{mainvectorspace}
		Let $3<p$ be a prime and $k$ a positive integer. Then we have
		\begin{enumerate}[\rm(1)] \vspace{-0.2cm}
			\item \[\sqrt{2/3}\cdot p^{k}<\sat(3-\mathrm{AP}, \F_p^{2k})\leq \left(\frac43+\frac{r}{3\cdot o_{p}(-2)}\right)(p^k-1),\]
			where $r$ is the residue modulo $3$ of the order $o_{p}(-2)$ of $-2$ in $\mathbb{F}_{p}^{\times}$.\vspace{-0.2cm}
			\item
			\[\sqrt{2/3}\cdot p^{k+\frac12}<\sat(3-\mathrm{AP}, \F_p^{2k+1})\leq \frac23(p^{k+1}+p^k-2)+\frac{r(p^{k+1}+p^k-2)}{3\cdot o_{p}(-2)},\]  %\todo[inline]{esetleg min(valséges becslés ; ez) Jo kérdés. Az a kérdés, hogy mit tekintunk jonak, $p$-t rogzitjuk és $k$-val tartunk a végtelenbe, vagy $k$-t rogzitjuk és $p$-vel tartunk a vegtelenbe. Utobbi esetben a legutolso, $(2,-1)$-szaturalo peldaink direktszorzatai is jobbat adnak mint a véletlen. (B)}
			where $r$ is the residue modulo $3$ of the order $o_{p}(-2)$ of $-2$ in $\mathbb{F}_{p}^{\times}$.
			\vspace{-0.2cm}
			\item 
			\[ p^{k}<\sat((2, -1), \F_p^{2k})\leq \left(\frac32+\frac{r}{2\cdot o_{p}(-2)}\right)(p^k-1),\]  
			where $r$ is the residue modulo $2$ of the order $o_{p}(-2)$ of $-2$ in $\mathbb{F}_{p}^{\times}$.\vspace{-0.2cm}
			
			%\textcolor{blue}{két "egyenes"-konstrukciós ps rend., plusz Remark 1.9. }
			\vspace{-0.2cm}
			\item \[ p^{k+\frac12}<\sat((2, -1), \F_p^{2k+1})\leq \sqrt{c_p}\left(\frac32+\frac{r}{2\cdot o_{p}(-2)}\right)(p^k-1)\sqrt{p},\]  
			where $r$ is the residue modulo $2$ of the order $o_{p}(-2)$ of $-2$ in $\mathbb{F}_{p}^{\times}$ and $c_p\leq 3$ is  the same constant depending on $p$ as in Theorem \ref{maingroups}.
			
			%\textcolor{blue}{két "egyenes"-konstrukció a $Z_m$-mel ötvözve direkt szorzattal, $(2,-1)$-gyel, ptn rend, plusz Remark 1.9. }
		\end{enumerate}
		%\todo{alsot is beírni?}
	\end{theorem}
	
	\begin{remark}
		The same ideas as in Theorem \ref{mainvectorspace} work to prove analogous results in $\mathbb{Z}_m^k$, when $\gcd(m,6)=1$. Then  $o_p(-2)$ should be replaced by the multiplicative order of $-2$ in the ring $\mathbb{Z}_m$. In the proofs Theorems \ref{dirprod} and \ref{dirprod2} should be used intead of Propositions \ref{linesp} and \ref{linesp2}, respectively. 
	\end{remark}
	
	\begin{theorem}\label{proby}
		For abelian groups $G$ of order $n > 5$ odd, it holds that
		$$\sat(3-\mathrm{AP},G) \leq \sat((1/2,1/2),G) \leq \sqrt{(n-1)\ln{(n-1)}}+ \sqrt{(n-1)}+1.$$    
	\end{theorem}
	
	
	
	%Our results use various techniques from  additive combinatorial and probabilistic tools to  finite geometric and design theoretic results.
	
	The paper is organized as follows. 
	In Section 2 we present some preliminary results and useful tools. First we prove lower bounds on the size of  saturating and complete $W$-avoiding sets. Then we show the strength and limitation of direct product constructions, which will enable us to prove the upper bound results for vector spaces (Theorems \ref{main3AP} and \ref{mainvectorspace}). To have upper bounds close to our lower bounds, we will rely on  further avoiding and saturating set constructions in finite fields which are small enough.  Finally, we point out the relation of these results to results concerning caps and $3$-AP covering sequences.\\
	In Section 3 we prove the upper bounds of Theorem \ref{main3AP} by analysing point sets of conics in the respected vector spaces.
	We also prove some general constructions of saturating sets in direct product of groups. Then we deduce the upper bounds of Theorem \ref{mainvectorspace} and Theorem \ref{main3AP} (1).
	Section 4 is devoted to the proof of Theorem \ref{maingroups} and \ref{proby}, where we provide constructions based on numeral systems, additive basis, Sidon sets and random constructions, using tools from additive number theory to design theory. 
	%A \textbf{ cap} in an affine or in a projective space  over a finite field $\F_q$ is  a  set  of  points  containing no collinear triples.  A  cap  is  called  \textbf{complete}  when  it  cannot  be  extended  to  a  large  cap.  The  central  problem  in  the  theory  of  caps  is  to  find  the  maximal  and  minimal sizes of caps in the affine or  projective space of given dimension, in terms of the order $q$ of the field.
	%As an application of our results on $\beta$-ratio-free sets, we show complete complete caps in $\AG(n,q)$ and point out that  
	
	\section{Preliminary results}
	
	\subsection{Double counting and direct sum constructions}
	
	
	
	We begin this section by demonstrating some trivial lower bounds on the size of $3$-$\mathrm{AP}$ saturating and $W$-saturating sets.
	
	\begin{prop}
		\label{satbound}
		\begin{enumerate}[\rm(1)]
			
			\item Suppose that $H$ is a $3$-$\mathrm{AP}$ saturating set in the abelian group $G$ of odd order. Then
			\[|H|\geq \sqrt{\frac{2}{3}|G|+\frac{1}{36}}+\frac{1}{6}.\]
			Hence, $\sat(3-\mathrm{AP}, \F_q^n)>0.8164\cdot q^{n/2}$.  \vspace{-0.2cm}
			\item Suppose that $H$ is a $w$-saturating set in the group $G$, where 
			$w\in \mathbb{Z}\times \mathbb{Z}$ %\{(2,-1),(1,-1)\}$, 
			or $w=(\lambda_1,\lambda_2)\in \F_q^*\times \F_q^*$  if $G=\F_q^n$. Then 
			\[|H|\geq \lceil \sqrt{|G|} \rceil.\]
			Hence, $\sat(w, \F_q^n)\geq \lceil q^{n/2} \rceil$.  \vspace{-0.2cm}
			\item Suppose that $H$ is a $w$-saturating set in the group $G$, where 
			$w=(1,1)$, or $w=(\frac12, \frac12)$ if $G$ has odd order, or $w=(\lambda,\lambda)$ for some $\lambda \in \F_q^*$ if $G=\F_q^n$. Then
			\[|H|\geq \sqrt{2|G|+\frac14}  -\frac12.\]
			So in this case $\sat(w, \F_q^n)>\sqrt2\cdot q^{n/2}-0.5$
		\end{enumerate}
	\end{prop}
	\begin{proof}
		Part $(1)$. 
		By definition, $\forall x\in G\setminus H$, we have $a\neq b\in H$ s.t.  $x=\frac12 a + \frac12 b$, or $x=2a-b$, or $x=-a+2b$. Thus by double counting, \begin{equation*}\label{eq:1}
			|G\setminus H|\leq 3\binom{|H|}{2},
		\end{equation*}
		from which the lower bound follows.
		
		Part $(2)$.
		Let $w=(\lambda_1, \lambda_2)$. 
		By definition, $\forall x\in G\setminus H$, we have $h\neq h'\in H$ s.t.  $x= \lambda_1h+\lambda_2h'$. 
		Then by double counting, \begin{equation*}\label{eq:11}
			|G|- |H|\leq 2\binom{|H|}{2}.
		\end{equation*}
		After rearranging, we get the desired bound.
		
		Part $(3)$. By double counting,
		\[|G|-|H|\leq \binom{|H|}{2},\]
		and the bound follows after rearranging.
	\end{proof}
	
	%\begin{prop}
	
	%\end{prop}
	
	%\begin{prop}\label{lower bound1}\todo{Ez a étel igy nem igaz, lasd. email.}
	%Suppose that $H$ is a $W$-saturating set in the Abelian group $G$, where $W$ contains a vector $(\lambda_1, \lambda_2)$ such that $\lambda_1\neq \lambda_2$, $\lambda_1+ \lambda_2=1$. Then we have $$|H|\geq \sqrt{|G|}.$$
	%\end{prop}
	
	%\begin{proof}
	%   By definition, $\forall x\in G\setminus H$, we have $h\neq h'\in H$   (Note that $h \neq h'$ must hold since $\lambda_1+ \lambda_2=1$ )  s.t.  $x= \lambda_1h+\lambda_2h'$. Thus by double counting, \begin{equation*}\label{eq:1}
	%      |G\setminus H|\leq 2\binom{|H|}{2}.
	% \end{equation*}
	%After rearranging, we get the desired bound.
	%   Now suppose that $|H|= \sqrt{|G|}-c$ for some positive $c>0$. Then applying the product construction of Proposition \ref{product} $t$ times, we get a set $H^t$ saturating $3$-$\mathrm{AP}$s in $G^t$. Note that this set has cardinality $(\sqrt{|G|}-c)^t$. However, using equation (\ref{eq:1}), we must have 
	%\begin{equation}\label{eq:2}
	%       |G|^t\leq 3\binom{(\sqrt{|G|}-c)^t}{2}+(\sqrt{|G|}-c)^t< 1.5(\sqrt{|G|}-c)^{2t}.
	%  \end{equation}
	% Choosing $t$ large enough such that $$\frac{|G|}{(\sqrt{|G|}-c)^2}>\sqrt[t]{1.5}$$ holds, this in turn leads to a contradiction.
	%\end{proof}
	
	
	% \begin{prop}
	% \label{satbound}
	%  %is at least roughly $\sqrt{2} \cdot q^{n/2}$.
	% %\todo[inline]{B: Erre lehet, hogy van jobb otleted, hogy lehetne szebben irni.\\Z: esetleg $sat>sqrt(2)*q^{n/2}-0.5$? nem éles mindig, de egzakt.} 
	% \end{prop}
	% \begin{proof}
	
	% \end{proof}
	
	\begin{remark}
		We will see later that the lower bound above for $\sat((2,-1),G)$ is sharp in some cyclic groups, cf. Theorems \ref{diffset} and \ref{gyok3}. Subsection \ref{last} provides further instances when Proposition \ref{satbound}  (2) is sharp. 
	\end{remark}
	
	%\begin{cor}
	%   $\sat(w, q^n)\geq q^{n/2}$ provided that $w=(\lambda_1, \lambda_2)$ s.t. $\lambda_1+\lambda_2=1$, $\lambda_1\neq\lambda_2$, and $\lambda_i\in \F_q^*$ for $i\in \{1,2\}$.
	%\end{cor}
	
	%\begin{remark} If we weaken the condition in Proposition \ref{lower bound1} and only require a vector  $(\lambda_1, \lambda_2)$ such that $\lambda_1\neq \lambda_2$, we get a slightly weaker \begin{equation*}\label{eq:1}
	%       |G\setminus H|\leq 2\binom{|H|}{2}+|H|.
	%  \end{equation*}
	%\end{remark}
	%The following proposition can be proved via a similar double counting, taking into account the $W$-saturation reformulation of $3$-$\mathrm{AP}$ saturating sets with .  $W=\{(2,-1), (1/2, 1/2)\}$. 
	%\begin{prop}
	%\label{lower bound2}
	%Suppose that the $H$ is a  $3$-$\mathrm{AP}$ saturating set in the Abelian group $G$. Then we have %$$|H|\geq \sqrt{\frac{2}{3}|G|+\frac{1}{36}}+\frac{1}{6}.$$
	%Hence, $\sat(3-\mathrm{AP}, q^n)>0.8164\cdot q^{n/2}.$
	%\end{prop}
	
	
	\noindent Next we show that the direct sum construction preserves certain properties concerning saturation and $W$-avoidance. Note however that saturation with respect to $3$-$\mathrm{AP}$s is not preserved.
	
	
	\begin{prop}[Avoiding and saturation property in direct products]\label{product}
		\ \\ \vspace{-0.5cm}
		\begin{enumerate}[\rm(1)]
			\item Suppose that $H$ and $H'$ are subsets of the abelian groups $G$ and $G'$, respectively.
			\begin{enumerate}
				\item 
				Assume that $H$ and $H'$ are $3$-$\mathrm{AP}$ free in the corresponding groups. Also, if $G$ (or $G'$) has even order, then assume  that the order of $x-y$ is larger than $2$ for any two distinct $x,y\in H$ ($\in H'$). 
				Then $H\times H'$ is $3$-$\mathrm{AP}$ free in $G \times G'$.
				\item If $H$ and $H'$ are $(1,1)$-avoiding in the corresponding groups and $0_G \notin H$, $0_{G'}\notin H'$ then $H\times H'$ is $(1,1)$-avoiding in $H \times H'$.
			\end{enumerate}
			
			%\todo{ez milyen megszorítással igaz általános W-free esetben? a nontrivial megoldás mit jelentsen pon tosan? Legjobb lenne úgy definiálni, hogy a 3Ap essen ki egyból mint kételemű vektorhalmaz következmény}
			
			\item  Suppose that $H$ and $H'$ are $W$-avoiding subsets of the vector spaces $\F_q^m$ and $\F_q^n$, respectively for some $W \subseteq \F_q^*\times \F_q^*$.  Then $H\times H'$ is $W$-avoiding in $\F_q^m \times \F_q^n$, provided that $\lambda_1+\lambda_2=1$ for all $w=(\lambda_1, \lambda_2)\in W$.
			
			\item Suppose that the set $W$  consists of a single vector $w=(\lambda_1, \lambda_2) \in \F_q^*\times \F_q^*$, provided that $\lambda_1+\lambda_2=1$. If $H\subseteq \F_q^m$ and $H'\subseteq \F_q^n$ are $W$-saturating sets of the corresponding vector space, then $H\times H'$ is also a $W$-saturating set in $ \F_q^{m}\times \F_q^n$.
			%  \item If $H$ and $H'$ are complete $3$-$\mathrm{AP}$ free in the corresponding groups, then $H\times H'$ is complete $3$-$\mathrm{AP}$ free in $G\times G'$.
			\item Suppose that $w=(2,-1)$, or $G$ and $G'$ are of odd order and $w=(1/2,1/2)$. If $H\subseteq G$ and $H'\subseteq G'$ are $w$-saturating sets of the corresponding groups, then $H\times H'$ is also a $w$-saturating set in $G \times G'$.
		\end{enumerate}
	\end{prop} %\todo{Bele kéne irni, hogy a $\lambda_1+\lambda_2=1$ feltételbol kovetkezik, hogy egy egyenesre esnek a vektorterben.}
	
	Note that the condition $\lambda_1+\lambda_2=1$ holds if and only if  $w_1$, $w_2 $ and   $\lambda_1w_1 + \lambda_2w_2 \in \F_q^n$ are collinear in the affine space $\F_q^n$.
	% As a consequence, by taking $W=\{(2,-1), (1/2, 1/2) \}$, part $(2)$ implies the $3$-$\mathrm{AP}$ freeness of the product structures in vector spaces.
	Observe also that by choosing $w=(2, -1)$ in part (3), 
	the direct product will be $3$-$\mathrm{AP}$ saturating as well.
	
	\begin{proof}
		%    The first part is (probably) a folklore. We prove it for the sake of completeness. 
		$(1a)$ Assume to the contrary that there is a $3$-$\mathrm{AP}$: 
		\[(h_1,h_1'),(h_2,h_2'),(h_3,h_3') \in  H\times H'\] 
		with difference $(d,d')\in G\times G'$. Since $(d,d')$ is not the neutral element of the group $G\times G'$, w.l.o.g. we may assume that $d$ is not the neutral element of $G$.
		Then $h_1,h_2,h_3$ is a $3$-$\mathrm{AP}$ of $G$, contradicting the assumption on $H$, or the order of $G$ is even and $h_1+h_3=2h_2$ holds because the size of $\{h_1,h_2,h_3\}$ is $2$, i.e. $h_1=h_3$ and hence the order of $h_1-h_2$ is $2$, a contradiction.
		
		$(1b)$ Assume to the contrary 
		$(h_1,h_1')=(h_2,h_2')+(h_3,h_3')$ for some $h_1,h_2,h_3 \in H$ and $h_1',h_2',h_3' \in H'$. W.l.o.g. we may assume $h_2 \neq h_3$. Then $h_1=h_2+h_3$ are $3$ distinct elements of $H$ contradicting the fact that $H$ is $(1,1)$-avoiding (recall $0_G\notin H$). 
		
		Proof of $(2)$. Assume to the contrary the existence of $w=(\lambda_1, \lambda_2 )\in W$ such that  $(h_1,h_1')= \lambda_1(h_2,h_2')+ \lambda_2(h_3,h_3')$ for some elements of $H\times H'$.  Hence
		\begin{equation*}   
			% \begin{enumerate}
			%   \item 
			\begin{cases}
				h_1=\lambda_1h_2+\lambda_2h_3,\\
				h_1'=\lambda_1h_2'+\lambda_2h_3'.\\
			\end{cases}
			%  \end{enumerate} 
		\end{equation*}
		Since $(h_2,h_2')\neq (h_3,h_3')$, we may assume w.l.o.g. that $h_2 \neq h_3$.
		We have $h_1=\lambda_1 h_2 + \lambda_2 h_3$ and this is a contradiction if $h_1, h_2, h_3$ are three pairwise distinct elements since $H$ is $W$-avoiding. If we had $h_1=h_2$, then $h_1(1-\lambda_1)=h_3(1-\lambda_1)$ and hence also $h_1=h_3$, a contradiction since $h_2 \neq h_3$ (and the same argument shows $h_1\neq h_3$ as well). %\todo{atirtam ezt a reszt}
		%Since $(h_1,h_1')\neq (h_2,h_2')$ and $(h_1,h_1')\neq (h_3,h_3')$, we may assume that $h_1\neq h_2$ w.l.o.g. Then $h_1\neq h_3$ holds as well, since $\lambda_1+\lambda_2=1$. This would yield a nontrivial solution $h_1, h_2, h_3$ for the equality $h_1=\lambda_1h_2+\lambda_2h_3$ in $H$, a contradiction.
		
		
		Proof of $(3)$. Take any $(g,g') \in (\F_q^{m}\times \F_q^{n}) \setminus (H \times H')$. If $g\notin H$ and $g'\notin H'$, then by the assumption, there exist $h_1,h_2 \in H$ and $h_1',h_2' \in H'$ such that
		\begin{equation*}   
			% \begin{enumerate}
			%   \item 
			\begin{cases}
				g=\lambda_1h_1+\lambda_2h_2,\\
				g'=\lambda_1h_1'+\lambda_2h_2',\\
			\end{cases}
			%  \end{enumerate} 
		\end{equation*}
		$g,h_1,h_2$ and $g',h_1',h_2'$ are sets of pairwise distinct elements.
		%where the choice $h_1=h_2$ or $h_1'=h_2'$ is allowed. 
		This in turn shows that $$(g,g')=\lambda_1(h_1, h_1')+\lambda_1(h_2, h_2'),$$
		where $(g,g'),(h_1,h_1'),(h_2,h_2')$ are pairwise distinct elements. 
		
		We cannot have $g\in H$ and $g'\in H'$ at the same time, hence w.l.o.g. we may assume $g\in H$ and $g'\notin H'$. 
		Then by the assumption, there exist $h_1',h_2' \in H'$ such that
		\[g'=\lambda_1h_1'+\lambda_2h_2',\]
		$g', h_1', h_2'$ are pairwise distinct elements.
		This in turn shows that $$(g,g')=\lambda_1(g, h_1')+\lambda_2(g, h_2'),$$
		where $(g,g'),(g,h_1'),(g,h_2')$ are pairwise distinct elements. 
		
		The proof of $(4)$ is the same as the proof of $(3)$.\end{proof}
	
	A generalisation of some of the results above will be discussed in Subsection \ref{last}.
	%\todo{Ez is uj.} 
	
	
	%If $g\notin H$ and $g'\notin H'$ then there exist $h_1,h_2 \in H$ and $h_1',h_2' \in H'$ such that $h_1,h_2,g$ and $h_1',h_2',g'$ are $3$-$\mathrm{AP}$s with differences $d\in H$ and $d'\in H'$, respectively. Then $(h_1,h_1'),(h_2,h_2'),(g,g')$ is a $3$-$\mathrm{AP}$ with difference $(d,d')$.
	% Otherwise w.l.o.g. we may assume $g\notin H$ and $g'\in H'$. Then there exist $h_1,h_2 \in H$ such that $h_1,h_2,g$ is a $3$-$\mathrm{AP}$ with difference $d$. Then
	%$(h_1,g'),(h_2,g'),(g,g')$ is a $3$-$\mathrm{AP}$ with difference $(d,0)$ and $(h_1,g'),(h_2,g')\in H \times H'$.
	%The third part is a direct consequence of the first two parts.
	
	
	% \begin{prop}\label{product}
	%     Suppose that $H$ and $H'$ are (not necessarily distinct)  complete $3$-$\mathrm{AP}$ free sets in a commutative group $G$. Then $H\times H'$ is a complete $3$-$\mathrm{AP}$ free set in $G\times G$
	% \end{prop}
	
	%ctr+shift+7 (for comments)
	
	% In the rest of the paper we will assume that $G$ has odd order.
	
	\begin{corollary}
		Direct product of complete $(2,-1)$-avoiding sets is a complete $3$-$\mathrm{AP}$ free set. Moreover,
		$  a((2,-1),G)\cdot a((2,-1),H) \ge a(3-AP,G\times H)$. This highlights the importance of finding complete $(2,-1)$-avoiding sets $A$ in $G$ such that  $|A|\leq \sqrt{|G|} $. 
	\end{corollary}
	
	%\todo[inline]{Tulajdonképpen a szorzat konstrukci\'oban az a j\'o amikor tényleg $|A|\leq \sqrt{|G|}$ (teh\'at a fels\H o egészrész nélk\"ul). Lehet, hogy a 'kicsit' fels\H o egéSzréSz nélkul kéne definialni. Ez esetben viszont a Singer példa nem lesz kicsi.}
	
	
	The previous propositions %\ref{lower bound1}, \ref{lower bound2} and \ref{product}
	motivate the distinguishment below.
	
	\begin{definition}
		A complete $3$-$\mathrm{AP}$-avoiding or a complete $(\lambda, 1-\lambda)$-avoiding set $H\subseteq G$ is called \textbf{small} if $|H|\leq  \sqrt{|G|} $. 
	\end{definition}
	
	
	\begin{proposition} 
		\label{mirejo}
		Let $H$ denote a $W$-avoiding, $W'$-saturating set in $\F_q^m$.
		\begin{enumerate}[\rm(1)]
			\item Then for each $\lambda \in \F_q^*$ it holds that $\lambda H$ is $W$-avoiding and $W'$-saturating.
			\item If for each $(\lambda_1,\lambda_2) \in W$ it holds that $\lambda_1+\lambda_2=1$, then for each $d\in \F_q^m$, $H + d$ is $W$-avoiding. If for each $(\lambda_1,\lambda_2) \in W'$ it holds that $\lambda_1+\lambda_2=1$, then for each $d\in \F_q^m$, $H + d$ is $W'$-saturating.
			\item Assume $W'=\{(1,1)\}$, $0 \in H$ and $\lambda \in \F_{q}^m \setminus \{0,1\}$. Then for each $x\in \F_q^m \setminus \{0\}$ there exist $a,b \in \lambda H$ such that $x=(1/\lambda) a + (1/\lambda) b$. In particular, $\lambda H$ is $(1/\lambda,1/\lambda)$-saturating.
		\end{enumerate}
	\end{proposition}
	\begin{proof}
		Proof of (1). For some $(\lambda_1,\lambda_2)\in W$ and $a,b,c \in H$, $\lambda_1 \lambda a + \lambda_2 \lambda b = \lambda c $ would imply $\lambda_1 a + \lambda_2 b = c$, a contradiction, which proves that $\lambda H$ is $W$-avoiding. Also, if $x\in \F_q^m \setminus \lambda H$, then $x=\lambda c$ for some $c \notin H$ and hence $c=\lambda_1 a + \lambda_2 b$ for some $(\lambda_1,\lambda_2) \in W'$. It follows that $\lambda H$ is $W'$-saturating.
		
		Proof of (2). If we had $\lambda_1(a+d)+\lambda_2(b+d)=c+d$ for some $a,b,c \in H$ and $(\lambda_1,\lambda_2)\in W$, then also $\lambda_! a+\lambda_2 b=c$, a contradiction. If $x \notin H+d$ then $x=c+d$ for some $c\notin H$ and hence there exists $(\lambda_1,\lambda_2)\in W'$ such that $\lambda_1 a + \lambda_2 b = c$ proving that $H+d$ is $W'$-saturating.
		
		Proof of (3). Take some $x\neq 0$.
		If $x \notin H$, then $x=a+b$ for some $a,b \in H$ and hence $x=(1/\lambda) (\lambda a)+(1/\lambda) (\lambda b)$. If $x \in H$, then $x=(1/\lambda) (\lambda 0)+(1/\lambda) (\lambda x)$. 
	\end{proof}
	
	
	
	% \begin{proposition}
	% For some $(\lambda_1,\lambda_2)\in (\F_q^*)^2$, $\lambda_1+\lambda_2\neq 0$, let $H \subseteq \F_q^m$ be such that for each $x\in \F_q^m$ there exist $a, b \in H$ such that $x=\lambda_1 a + \lambda_2 b$. Then 
	% \[\frac{1}{\lambda_1+\lambda_2}((\lambda_1+\lambda_2)H\times (\lambda_1+\lambda_2)H \times \ldots \times (\lambda_1+\lambda_2)H)\]
	% is $(\lambda_1,\lambda_2)$-saturating in $\F_q^m \times \F_q^m \times \ldots \times \F_q^m$.
	% \end{proposition}
	% \begin{proof}
	%     For each $x\in \F_q^m$ there exist $a,b \in (\lambda_1 + \lambda_2)H$ such that 
	%     \[x=\frac{\lambda_1}{\lambda_1+\lambda_2} a + \frac{\lambda_2}{\lambda_1+\lambda_2}b\]
	% \end{proof}
	
	\begin{proposition}
		Let $q$ be odd. If $H \subseteq \F_q^m$ and $H' \subseteq \F_q^n$ are $(1,1)$-saturating such that the corresponding zero vectors are contained 
		in $H$ and in $H'$, resp.,  then $\frac12(2H \times 2H')$ is $(1,1)$-saturating in $\F_q^m \times \F_q^n$. 
	\end{proposition}
	\begin{proof}
		By (3) of Proposition \ref{mirejo} it follows that $2H$ is $(1/2,1/2)$-saturating in $\F_q^m$ and the same holds in $\F_q^n$ for $2H'$. Then by (3) of Proposition \ref{product} it follows that $(2H)\times (2H')$ is $(1/2,1/2)$-saturating in $\F_q^m \times \F_q^n$. Since this subset contains the zero vector of $\F_q^m \times \F_q^n$, the statement follows again by (3) of Proposition \ref{mirejo}. 
	\end{proof}
	
	\begin{proposition}
		If $H \subseteq \F_q^m$ and $H' \subseteq \F_q^n$ are $(1,-1)$-saturating such that the corresponding zero vectors are contained in $H$ and in $H'$, then $H \times H'$ is $(1,-1)$-saturating in $\F_q^m \times \F_q^n$.  
	\end{proposition}
	\begin{proof}
		Take any $(g,g') \in (\F_q^{m}\times \F_q^{n}) \setminus (H \times H')$. If $g\notin H$ and $g'\notin H'$, then by the assumption, there exist $h_1,h_2 \in H$ and $h_1',h_2' \in H'$ such that
		\begin{equation*}   
			% \begin{enumerate}
			%   \item 
			\begin{cases}
				g=h_1-h_2,\\
				g'=h_1'-h_2',\\
			\end{cases}
			%  \end{enumerate} 
		\end{equation*}
		$g,h_1,h_2$ and $g',h_1',h_2'$ are sets of pairwise distinct elements.
		%where the choice $h_1=h_2$ or $h_1'=h_2'$ is allowed. 
		This in turn shows that $$(g,g')=(h_1, h_1')-(h_2, h_2'),$$
		where $(g,g'),(h_1,h_1'),(h_2,h_2')$ are pairwise distinct elements. 
		
		We cannot have $g\in H$ and $g'\in H'$ at the same time, hence w.l.o.g. we may assume $g\in H$ and $g'\notin H'$. 
		Then by the assumption, there exist $h_1',h_2' \in H'$ such that
		\[g'=h_1'-h_2',\]
		$g', h_1', h_2'$ are pairwise distinct elements.
		This in turn shows that $$(g,g')=(g, h_1')-(0, h_2'),$$
		where $(g,g'),(g,h_1'),(0,h_2')$ are pairwise distinct elements. 
	\end{proof}
	
	
	%\subsection{More on the direct sum construction}
	
	\subsection{Relation to caps}
	
	Let $q$ denote  {\em any} (even or odd) prime power. We describe the relation of the results above  to results concerning caps.
	
	\begin{definition}
		A cap of $\AG(n,q)$ is a point set meeting each line of $\AG(n,q)$ in at most two points. 
		
		\smallskip
		
		A cap is called complete if it cannot be extended to a larger cap.
		
		\smallskip
		
		A saturating set $S$ of $\AG(n,q)$ is a point set with the property that for each $P \in \AG(n,q) ~\setminus~ S$ there exist two distinct points $Q,R \in S$, such that $P$ is incident with the line joining $Q$ and $R$.
	\end{definition}
	
	It is clear from the definitions above that a cap is complete if and only if it is also a saturating set.\\
	The lattice of affine subspaces of $\F_q^n$ is isomorphic to the subspace lattice of $\AG(n,q)$. %The lines of $\AG(n,q)$ correspond to the one-dimensional linear varieties of $\F_q^n$, thus we may use the definitions above for subsets of $\F_q^n$ as well.\\
	For results on the (maximum) size of complete caps in $\AG(n,q)$, we refer to \cite{EB, Edel, EPPP, Ty} and the references therein. 
	%\todo[inline]{Ezek nem végtelen sok $q$-ra mukodo konstrukciok. Kulon szedni a $q=3$ illetve $Z_m$-es cikkeket.}
	A related problem is the smallest size of complete caps in finite affine and projective spaces. For the size of small complete caps, the theoretical lower bound is essentially sharp for $q$ even \cite{evencap4, evencap3, evencap2, evencap1} and in some cases also for $q$ odd \cite{completecap}. See also \cite{Anbar, Bartoli, DO, Giuli} and the references therein for small complete caps for $q$ odd.
	%Results concerning the smallest size of complete caps in $\AG(n,q)$,  can be find in \cite{Giulietti, Kara, Oster, Platon}. Anbar 6.2-es tétele itt a lényeg!
	
	
	
	\begin{prop}
		Put $W=\{(\lambda_1, \lambda_2) \in \F_q^*\times \F_q^* : \lambda_1+\lambda_2=1\}$. Then we obtain the following.
		\begin{enumerate}[\rm(1)]
			\item Caps of $\AG(n,q)$ and $W$-avoiding sets of $\F_q^n$ are equivalent objects. In particular, by Theorem \ref{product} Part $(2)$ the direct sum of caps is a cap.  
			\item Saturating sets of $\AG(n,q)$ and $W$-saturating sets of $\F_q^n$ are equivalent objects.
			\item Complete caps of $\AG(n,q)$ and complete $W$-avoiding sets of $\F_q^n$ are equivalent objects.
		\end{enumerate}
		If $q=3$, then $W=\{(2,2)\}$, if $q=4$, then 
		$W=\{(i,1+i)\}$. Hence, by Part $(3)$ of Theorem \ref{product}, direct sum of saturating sets is a saturating set and direct sum of complete caps is a complete cap for $q\in \{3,4\}$.  \qed
	\end{prop}
	
	
	
	\subsection{Related results on solving linear equations in algebraic structures}
	
	We summarize some results which are connected to the theme of this paper in the sense that the subject is a subset of a set, in which an equation of special form has no non-trivial solutions. Then we also mention some results of saturation type with respect to an equation.%when every element outside the subset are determined via the equation (in some sense) using the elements of the subset.
	%\todo{Ezt a mondatot nem annyira ertem itt.}
	
	Most probably the leading examples for the first theme are the Sidon sets.
	A set of elements in an abelian group is called a Sidon set if all pairwise sums of its not necessarily distinct elements are distinct. Equivalently, the equation 
	$a+b=c+d$  has only the trivial solution 
	$\{a, b\}= \{c,d\}$ in the set. Observe that Sidon sets are $3$-$\mathrm{AP}$ free.
	Concerning Sidon sets, Erdős and Turán  observed  \cite{ET} (see also Cilleruelo  \cite{Cille}) that the point set of the parabola in $\F_q\times \F_q$ provides a Sidon set. Cilleruelo showed several further abelian groups admitting Sidon sets of size equal roughly to the square root of the order of the group.  Building on his observations, Huang, Tait and Won showed \cite{Tait} that the largest Sidon sets in $\F_3^n$ are of size $3^{n/2}$, provided that $n$ is even.
	Small complete Sidon sets of abelian $2$-groups are investigated in the recent paper of G. Nagy \cite{GNagy} who showed constructions gained from elipses and hyperbolas in the finite affine plane $\F_q\times \F_q$, and in the papers \cite{Pott} and \cite{red}.
	
	%Note that if $0\in A$ is a Sidon set, then it is necessarily $(1,-1)$-avoiding and it is not hard to see that the point set of the parabola in $\F_q\times \F_q$ is complete $(1,-1)$-avoiding, thus 
	
	%\begin{prop}
	%$a((1,-1),\F_p^{2n})= p^n$ for every odd prime $p$.    \qed
	%\end{prop}
	%\begin{proof}
	%For each $(a,b)\in \F_q\times \F_q$, $b\neq a^2$, there exist $x,y \in \F_q$ such that $(x,x^2)-(y,y^2)=(a,b)$. Indeed, put $x=(b/a+a)/2$, $y=(b/a-a)/2$.
	%\end{proof}
	
	There is a strong connection between  the case of vector space or cyclic group setting 
	and the case of integer setting, when  a non-trivial solution of a particular equation is forbidden within the interval
	$[1,n]\subset \mathbb{Z}$. %\todo{Nem értem ezt a mondatot.}.
	For Sidon sets, the papers of Ruzsa \cite{Ruzsa2, Ruzsa} discuss the case of small complete structures, while the  work of Kiss, Sándor and Yang \cite{SandorCs} deals with small saturating sets with respect to $3$-$\mathrm{AP}$s. 
	%\todo[inline]{Atirni a konstanst, illete definialni a $3$-AP covering sequencet, mert mas mint nalunk.}
	They used the term \textit{$3$-$\mathrm{AP}$ covering} sequence for another related concept. 
	Let $A_0=\{a_1<\ldots<a_t\}$
	be a set of nonnegative integers such that $\{a_1<\ldots<a_t\}$
	does not contain a 3-term arithmetic progression. A sequence $A=\{a_1,a_2,\ldots\}$
	is called the Stanley sequence of order  $3$ generated by $A_0$, where  the elements outside $A_0$
	are defined by a greedy algorithm as follows. For any $l\ge t$, $a_{l+1}$  is the smallest integer $a>a_l$
	such that $\{a_1,\ldots,a_l\}\cup\{a\}$
	does not contain a $3$-term arithmetic progression. Moreover, a sequence $A$ of non-negative integers is called a $3$-$\mathrm{AP}$-covering sequence if there exists an integer $n_0$
	such that, if $n>n_0$, then there exist $a_1, a_2 \in A$ such that $a_1, a_2, n$
	form a $3$-term arithmetic progression.
	Fang \cite{Fang} made the following improvement.
	
	
	\begin{theorem}[Fang \cite{Fang}]\label{fange}
		There is a $3$-$\mathrm{AP}$ covering sequence $S$ of integers such that  $$ \frac{ \mid S \cap [1,n]\mid}{\sqrt{n}}\leq \frac{8}{\sqrt{5}} \approx 3.578$$ holds for all $n$. 
	\end{theorem}
	\noindent This constant cannot be improved to $1.77$ \cite{SandorCs}.
	Note that this  result is strongly related to our main problem since it in turn shows that  $\sat((2,-1), \mathbb{Z}_n )\leq (3.578+o(1))\sqrt{n}$.
	
	
	
	\section{\texorpdfstring{Constructions in $\F_q \times \F_q$ and in other direct products}{Construction in Fq x Fq and in other direct products}} %and in $\F_q^{2t}$}{Construction in Fq x Fq}}
	
	In this section $q$ always denotes a prime power, and $p$ denotes a prime.
	%\begin{remark}
	%   If $S$ saturates the $3$-$\mathrm{AP}$s of $\F_q\times \F_q$, then $S$ is a saturating set of $\AG(2,q)$. %by adding $2$ points of the line at infinity to $S$ we obtain a saturating set of $\PG(2,q)$
	%\end{remark}\todo{ez igaz, de ez nagy satuhalmazt ad; nem tudom emiatt érdekes-e.}
	
	
	\begin{constr}\label{parabola}
		Let $\mathcal{P}$ be the point set of the parabola 
		\[\{(x,x^2) : x \in \F_q\}.\]
		%defined by  $y~=~x^2 : x, y\in \F_q$.
	\end{constr}
	
	\begin{theorem}\label{main1} Let $q$ be an odd prime power.
		If $-2$ is not a square in $\F_q$ then Construction \ref{parabola} is a complete $3$-$\mathrm{AP}$ free subset of $\F_q\times \F_q$, hence $a(3-\mathrm{AP}, \F_q^2)\leq q.$
	\end{theorem}
	\begin{remark}
		Note that  Construction \ref{parabola} provides an infinite family of small complete $3$-$\mathrm{AP}$ free subsets of $\F_q\times \F_q$.
		As observed by Erdős and Turán, the parabola construction provides also a (dense) Sidon set, see \cite{Eber, ET}.
		%\todo{Ezt a megjegyzest totolni? Vagy atirni arra, hogy infinite family of small complete $3$-$\mathrm{AP}$ free sets...}   
	\end{remark} 
	
	%\begin{cor} 
	%For every even $t$, Theorem \ref{main1} implies the existence of an infinite family of complete $3$-$\mathrm{AP}$ free subsets of $\F_q^t$, which attain the corresponding lower bound $q^{t/2}$ of Theorem \ref{lower bound1}, provided that $-2$ is not a square element in $\F_q$. This follows from the application of the product construction of Proposition \ref{product}.
	%\end{cor}
	
	\begin{proof}[Proof of Theorem \ref{main1}]
		For each $(a,b)\in \F_q^2$, $b\neq a^2$, we prove that one of the following systems of equations have a solution $(x,y)\in \F_q\times \F_q$. 
		\begin{equation*}   
			% \begin{enumerate}
			%   \item 
			\begin{cases}
				x+y=2a \\
				x^2+y^2=2b\\
			\end{cases}
			% \item 
			\begin{cases}
				2y-x=a \\
				2y^2-x^2=b\\
			\end{cases}
			%  \end{enumerate} 
		\end{equation*}
		This implies that no point $(a,b)$ outside $\cP$ can be added to the construction without violating the  $3$-$\mathrm{AP}$ free property.\\
		Solution for the first system exists if and only if $b-a^2$ is a square in $\F_q$. Indeed, in order to have a common solution, we should get a square value for the discriminant $16a^2-8\cdot(4a^2-2b)=16(b-a^2)$ of  $x^2+(x^2-4ax+4a^2)-2b=0$.\\
		Solution for the second system exists if and only if $2(a^2-b)$ is a square in $\F_q$. Indeed, in order to have a common solution, we should get a square value for the discriminant $16a^2-8\cdot(a^2+b)=8(a^2-b)$ of  $2y^2-(4y^2-4ay+a^2)-b=0$.\\
		If $-2$ is not a square, then either the first, or the second discriminant will be a square, providing a solution to one of the systems.
	\end{proof}
	
	Theorem \ref{main1} in turn implies the upper bound of Theorem \ref{main3AP} (1) on complete $3$-AP free sets in vector spaces once one notes that $-2$ is not a square element in $\mathbb{F}_p$ if and only if it is not a square in $\mathbb{F}_p^{2k+1}$.
	
	Now we show some saturating set constructions.
	
	\begin{constr}
		\label{lines} Let $\langle -2 \rangle$ denote the multiplicative (cyclic) subgroup of $\F_q^\times$, generated by $-2$, where char$(q)\neq 2,3$. Take a set of maximum size in each coset of $\langle -2 \rangle$ for which the equations $-2g=g'$, $4g=g'$ have no solutions within the set. Let $R$ denote the union of these sets.\\
		Let $\cL$ denote the set of point $\cL=\{(0,r) :  r\in \F_q^* \setminus R  \}\cup\{(r,0): r\in \F_q^* \setminus R  \}\subset \F_q \times \F_q$ 
	\end{constr}
	
	The following result is a reformulation of the upper bound of Theorem \ref{mainvectorspace} (1).
	
	\begin{prop}
		\label{linesp}
		Construction \ref{lines} contains $$2(q-1)\frac{o_q(-2)-\lfloor o_q(-2)/3\rfloor}{o_q(-2)}$$ elements and it saturates the $3$-$\mathrm{AP}$s of $\F_q\times \F_q$.  
	\end{prop}
	\begin{cor}
		If $3 \mid o_q(-2)$ then $|\cL|=\frac{4}{3}(q-1)$.
	\end{cor}
	
	\begin{proof}[Proof (of Proposition \ref{linesp})]
		First,  observe that the choice of $R$ ensures that for all $g\in \F_q\setminus \{0\}$, at least two of $g, -2g, 4g$ are admissible coordinates in $\cL$. This implies that at most $$(q-1)\frac{\lfloor o_q(-2)/3\rfloor}{o_q(-2)}$$ elements are contained in $R$. On the other hand, choosing  each element of form $(-2)^{3t-1} $ such that $ 0<3t\leq o_q(-2)$ in $\langle -2 \rangle$ and applying similar rule in each coset yield equality in the bound above. Then the cardinality of points in $\cL$ follows.\\
		Then take any point $(a,b)\in \F_q\times \F_q$, where $a\ne 0\ne b$ and suppose that the addition of $(a,b)$ to the construction does not create a $3$-$\mathrm{AP}$.
		Now take $$(0,2b), (a,b), (2a,0);$$
		$$(-a, 0), (0, b/2), (a,b);$$ and 
		$$(0,-b), (a/2, 0), (a, b)$$ which form three disjoint $3$-$\mathrm{AP}$s consisting of two points of of the axes and $(a,b)$. Here we use the fact  that $o_q(-2)>2$.
		Since at most one element of $\{a/2, -a, 2a\}$ and of $\{b/2, -b, 2b\}$ is contained in $R$, $\cL$ will contain at least $4$ of the points listed above thus together with $(a,b)$, a $3$-$\mathrm{AP}$ would be formed, a contradiction.\\
		Finally, suppose that $a=0$ or $b=0$. Then the addition of $(a,b)$ would again provide at least one $3$-$\mathrm{AP}$, since $(a,b)$ would induce $\frac{q-1}{2}$ pairs $P, P'$ on the axis incident to $(a,b)$ for which $(a,b)$ is the midpoint of $P$ and $P'$, but the number of points on the axis in $\cL$ is larger than $q/2$ thus by the pigeon-hole principle, there would be a pair $P, P'\in \cL$ for which $(a,b)$ is a midpoint, hence the addition of $(a,b)$ is not allowed.
	\end{proof}
	
	\begin{remark}
		If $q$ is a power of the prime $p$ then the multiplicative order of $-2$ in $\F_q^{\times}$ is the same as the multiplicative order of $-2$ in $\F_p^{\times}$.
	\end{remark}
	
	Proposition \ref{linesp} implies directly the upper bound of Theorem \ref{mainvectorspace}  (1) in view of the previous remark if we apply  $q=p^{k}$. 
	
	To get the upper bound when the dimension of the vector space is odd (Theorem \ref{mainvectorspace} (4)), we modify the construction in a way that it $(2, -1)$-saturates the whole space. It enables us to apply the direct sum construction once we have a suitable general upper bound on $\sat((2, -1), \F_p)$.
	
	
	
	\begin{constr}
		\label{lines(2,-1)} Let $\langle -2 \rangle$ denote the multiplicative (cyclic) subgroup of $\F_q^\times$, generated by $-2$, where char$(q)\neq 2,3$. Take a set of maximum size in each coset of $\langle -2 \rangle$ for which the equations $-2g=g'$,  have no solutions within the set. Let $R^*$ denote the union of these sets.\\
		Let $\cL^*$ denote the set of points $\cL^*=\{(0,r) :  r\in \F_q^* \setminus R^*  \}\cup\{(r,0): r\in \F_q^*   \}$. 
	\end{constr}
	
	
	\begin{prop}
		\label{linesp2}
		Construction \ref{lines(2,-1)} contains $$2(q-1)-\frac{(q-1)\lfloor o_q(-2)/2\rfloor}{o_q(-2)}$$ elements and it is a $(2, -1)$-saturating  set (and hence $3-\mathrm{AP}$ saturating set)  of $\F_q\times \F_q$.  
	\end{prop}
	
	\begin{proof}
		One should observe that each point $(a,b)$, $a\neq 0\neq b$ is $(2,-1)$-saturated by the pairs  $(-a, 0), (0, b/2)$ and 
		$(0,-b), (a/2, 0)$, and at least one of these pairs will be contained in the construction. The cardinality of $\cL^*$ follows similarly to that of $\cL$ in the proof of Proposition \ref{linesp}.
	\end{proof}
	
	The result above in turn implies the upper bound of Theorem \ref{mainvectorspace} (3). Then,  Theorem \ref{mainvectorspace} (4) follows from the direct sum construction, Proposition \ref{product}, applying it to Construction \ref{lines(2,-1)} with $q=p^k$ and the $(2, -1)$-saturating set construction for $\F_p$, given in the next section (see also Theorem \ref{maingroups} (1)). 
	
	%\todo[inline]{nem merőleges egyenesekkel ugyanez a konst?}
	
	%\todo[inline]{What happens if the two lines are $y=0$ and $x=y$? Then the 3APs through $(a,b)$ are:  
	%\[(2b,2b), \quad (a,b), \quad (2a-2b,0);\]
	%\[(b-a,0),  \quad (b/2,b/2), \quad (a,b);\]
	%\[(a,b),  \quad (a/2-b/2,0), \quad (-b,-b);\]
	%It gives the same. I think that it will always depend on the order of $2$.}
	
	Along the same lines, one can prove the existence of $3$-AP saturating sets in direct products of abelian groups.
	
	Let $A$ and $B$ denote two abelian groups (written additively) of orders $a$ and $b$, respectively, such that $\gcd(ab,6)=1$.
	
	For  any element $g$ which is not the neutral element of the group, put $D_g=\{g,-2g,4g,-8g,\ldots,-g/2\}$. Since $\gcd(6,ab)=1$, the elements $g,-2g,4g,-8g$ are pairwise distinct, so $|D_g| \geq 4$.
	
	For any element $g$ which is not the neutral element of the group, we denote by $R_g$ a subset of $D_g$ of maximum size such that the equations $x=-2y$ and $x=4y$ cannot be solved within $D_g$. 
	Note that \[\frac13 |D_g| \geq |R_g|=\lfloor\frac13|D_g|\rfloor\geq \frac15 |D_g|.\] 
	Using these notation, we have 
	
	\begin{theorem}
		\label{dirprod}
		Let $A$ and $B$ denote two abelian groups (written additively) of orders $a$ and $b$, respectively, such that $\gcd(ab,6)=1$.
		Then
		\[\cL=\{(r,0_B) : r \neq 0_A,\, r\notin R_g \mbox{ for each } g \in A\} \cup \{(0_A,r) : r \neq 0_B,\, r\notin R_g \mbox{ for each } g \in B\}\]
		is $3$-$\mathrm{AP}$ saturating in $A\times B$.
		On the size of $\cL$ we have
		\[\frac43 (\sqrt{|A\times B|}-1)\leq \frac23(a+b-2) \leq |\cL|\leq \frac45(a+b-2) .\]
		\qed
	\end{theorem}
	
	Note that if $|D_g|$ is the same for each $g\in A$ and $g\in B$ where $g$ is different from the neutral element, then the size of $\cL$ can be expressed via $a$, $b$ and $|D_g|$ for a single $g$. %and the size of $A$ and $B$. 
	
	This leads to the statement of Theorem \ref{mainvectorspace} (2) by choosing $A=\F_{p}^{k+1}$ and $B=\F_p^k$.
	
	
	For any element $g$ of a group $A$ which is not the neutral element of the group, we define $D_g$ as before and we denote by $R^*_g$ a subset of $D_g$ of maximum size such that the equations $x=-2y$ cannot be solved within $D_g$. As before, the size of $D_g$ is at least $4$ and hence
	\[\frac12 |D_g| \geq |R^*_g| \geq \lfloor \frac12|D_g| \rfloor \geq \frac25|D_g|.\]
	Using these notation we have
	
	\begin{theorem}
		\label{dirprod2}
		Let $A$ and $B$ denote two abelian groups (written additively) of orders $a$ and $b$, respectively, such that $a$ is odd and $\gcd(b,6)=1$. 
		Then
		\[\cL^*=\{(r,0_B) : r \neq 0_A,\, r\in A\} \cup \{(0_A,r) : r \neq 0_B,\, r\notin R^*_g \mbox{ for each } g \in B\}\]
		is a $(2,-1)$-saturating (and hence $3-\mathrm{AP}$ saturating) set of $A\times B$.
		On the size of $\cL^*$ we have
		\[a+\frac 12 b- \frac32 \leq |\cL^*|\leq a+\frac35 b - \frac85.\]\end{theorem}
	\qed
	
	
	
	% \begin{corollary}
	
	% If $-2$ is not a square in $\F_q$ and $k=2\ell+1$, then $-2$ is not a square in $\F_{q^k}$, hence, according to Construction \ref{parabola}, $\{(x,x^2) : x \in \F_{q^k}\}$ is a complete $3$-$\mathrm{AP}$ free in $\F_{q^k}^2$. Since 
	% $\F_{q^k}^2\cong \F_q^{4\ell+2}$ as groups, it follows that in this case there exist complete $3$-$\mathrm{AP}$ free subsets of size $q^{2\ell+1}$ in $\F_q^{4\ell+2}$.
	
	% \end{corollary}
	
	
	\section{\texorpdfstring{Complete $3$-$\mathrm{AP}$ free sets and saturation in abelian groups}{Complete 3-AP free sets and saturation in abelian groups}}
	
	
	
	\subsection{Probabilistic upper bound on saturating sets}
	
	%\todo[inline]{Jo lenne adni valami peldat amikor ez a legjobb amit tudunk, talan $GF(3)^n$? Vagy elmagyarazni, hogy ez nem azert erdekes mert kicsit ad, hanem mert megmondjuk, hogy a véletlen mit ad - és esetleg igy atfogalmazni az allitast.\\ Válasz: Beírtam ide valamit, valóban fontos a magyarázat (Z)}
	
	We start with a general bound using probabilistic arguments and prove Theorem \ref{proby}. While it is off by a logarithmic factor from the lower bound, it is still the best we know in several cases (although not in vector spaces). It also highlights the algebraic nature of constructions meeting or being close to the lower bound.
	
	
	\begin{theorem}\label{random_satu} Suppose that the set $H$ saturates the $3$-$\mathrm{AP}$s in  the abelian group $G$ of order $n$, $n>5$ odd, and $H$ is of minimum size. Then we have $$|H|\leq \sqrt{(n-1)\ln{(n-1)}}+ \sqrt{(n-1)}+1.$$    
	\end{theorem}
	
	\begin{proof} The proof follows the probabilistic argument of \cite{Nagy} of the second author, on the size of saturating sets of projective planes.\\ Let $H_0$ be a random subset of $G$ 
		consisting of elements $g\in G$ where each element is chosen independently, uniformly at random with probability $p$. The parameter $p$ will be determined later on.
		Let $H_1$ be the set of elements $g\in G$ which  can be obtained as 
		$2g= h+h'$ or $g= 2h-h'$ for $h,h'\in H_0$.
		Let $X$ denote the random variable which takes the cardinality of $H_0$ and $Y $ denote the random variable which takes the cardinality of $H_1$. 
		
		Then $H_0\cup (G\setminus H_1)$ will provide a set $H$ that saturates the $3$-$\mathrm{AP}$s in $G$. We will determine the value of $p$ which minimise the expected value of $X+n-Y$.
		Clearly, $\E(X)=pn$. \\
		We call a pair $g_1, g_2$\textit{ induced by $g$} if $g_1+g_2=2g$. Hence each element of $G\setminus \{g\}$ is contained in exactly one pair induced by a fixed element $g$. If $g\not \in H_1$ then $H_0$ contains at most one element from each pair induced by $g$.
		Thus $$\pP(g\not \in H_1)<(1-p^2)^{\frac{1}{2}(n-1)}.$$
		By the linearity of expectation, we get 
		$$\E(X+n-Y)<n\left(p+(1-p^2)^{\frac{1}{2}(n-1)}\right).$$
		If $p=\sqrt{\frac{\ln{(n-1)}}{{n-1}}}$, this provides the existence of a set which saturates $3$-$\mathrm{AP}$s and have cardinality at most $$n\left(\sqrt{\frac{\ln{(n-1)}}{{n-1}}}+\left(1-\frac{\ln{(n-1)}}{n-1}\right)^{\frac{n-1}{2}}\right)<\sqrt{(n-1)\ln{(n-1)}}+1+\sqrt{n-1},$$
		taking into account that 
		
		$\sqrt{\frac{\ln{(n-1)}}{{n-1}}}+\frac{1}{\sqrt{n-1}}<1$ and applying the Bernoulli bound $(1-\frac{x}{m})^m<e^{x}$ for $x=-\ln{(n-1)}$ and $m=n-1$.
	\end{proof}
	Actually this argument shows that $\sat((1/2, 1/2), G)\leq  \sqrt{(n-1)\ln{(n-1)}}+ \sqrt{(n-1)}+~1$. The same bound can be easily obtained for $w=(2, -1)$-saturation as well.
	
	\begin{remark}
		Using the Lovász local lemma, one can prove that the probability of  $\pP(g\not \in~ H_1)$ can be  bounded from below by $(1-p^2)^{c(n-1)}$ for some positive constant $c$, which implies that the order of magnitude of a random construction obtained as above will be $\Theta(\sqrt{n \ln {n}})$.
	\end{remark}
	
	
	
	
	\subsection{\texorpdfstring{Complete $(2,-1)$-avoiding sets of minimum size in cyclic groups via difference sets}{Small complete (2,-1)-avoiding sets in cyclic groups via difference sets}}
	
	%\begin{constr}
	%\ Assume that $q=2^{2n}+2^n+1$ is a prime  and denote by $D'$ a (Singer)  $(q,2^n+1,1)$-difference set of the group $(\mathbb{Z}_q,+)$. 
	%\end{constr}
	
	The Singer difference sets of the cyclic group of order $q^2+q+1$, $q$ a prime power, provide maximal Sidon sets. These constructions inspire the construction below. We use without explicit reference the most well known facts concerning difference sets according to the Handbook of Combinatorial Designs \cite{handbook}.
	
	\begin{theorem}\label{diffset} 
		Put $M=2^{2n}+2^n+1$ and denote by $D'$ a Singer  $(M,2^n+~1,1)$-difference set of the cyclic group $(\mathbb{Z}_M,+)$. 
		Then $D'$ is a complete $3$-$\mathrm{AP}$ free subset of $\mathbb{Z}_M$. Moreover, $D'$ is complete $(2,-1)$-avoiding of size $\lceil \sqrt{M}\rceil$, so its size reaches the lower bound in Proposition \ref{satbound} part $(2)$. 
	\end{theorem}
	
	\begin{remark}
		Note that $M=2^{2n}+2^n+1$ is a prime for $n\in \{1,3,9\}$ and in these cases we obtain  complete $3$-$\mathrm{AP}$ free subsets in the corresponding finite fields of size $2^{2n}+2^n+1$. In general, $n$ needs to be a power of $3$ for this to hold.  Indeed, $M$ can be written as $M= \frac{2^{3n}-1}{2^n-1}$, and if there exists a proper divisor $d\mid 3n$ which is not a divisor of $n$, then $gcd(d,n)<d$. Now, by applying $gcd(2^d-1,2^n-1)=2^{gcd(d,n)}-1$, we get the identity \begin{equation}
			(2^{\gcd(d,n)}-1)\cdot r\cdot \frac{2^{3n}-1}{2^d-1}=(2^n-1)\cdot M,   
		\end{equation} where
		$2^d-1=r\cdot (2^{\gcd(d,n)}-1)$. Hence $r\mid M.$ But on the one hand,\\ $r\leq 2^d-1<2^{2n}<M$, on the other hand, $r>1$ as $\gcd(d,n)<d$.
	\end{remark}
	
	\begin{proof}[Proof of Theorem \ref{diffset}]
		According to the First Multiplier Theorem, for a translate $D$ of $D'$ it holds that $2D=D$. We will show that $D$ is complete $3$-$\mathrm{AP}$ free. Note that this implies that the translates of $D$ are complete $3$-$\mathrm{AP}$ free as well.
		
		% First of all, if $2D=D$, then $0\notin D$ since otherwise $a=a-0=2a-a$ for each $a\in D$ (and with $0,a,2a\in D$) contradicting the fact that $D$ is a difference set. 
		
		First we show that $2a=b+c$ cannot hold with $a,b,c\in D$ pairwise distinct elements. Indeed, it would imply $a-b=c-a$, contradicting the fact that $D$ is a difference set. It follows that $D$ is $3$-$\mathrm{AP}$ free.
		
		Since for each $a\in D$ we have also $2a \in D$, in $D \cup \{0\}$ we have the $3$-$\mathrm{AP}$: $\{0,a,2a\}$. This shows $0\notin D$ and that $D$ saturates $\{0\}$.\\    
		Now take any $g\in \mathbb{Z}_M \setminus D$, $g\neq 0$. Then there exist $a,b\in D$  such that $a-b=g$ and since $2D=D$, we have also $a=2c$ for some $c\in D$, that is, $2c=b+g$. We cannot have $c=b$ since in that case $c=g \in D$, a contradiction.
		
		The size of $D$ reaches the lower bound in Proposition \ref{satbound} (2) since $2^n<\sqrt{M}<2^n+1=|D|$.
	\end{proof}
	
	\begin{corollary}
		For $p\in \{7,73,262657\}$, the minimum size of a complete $(2,-1)$-avoiding subset of $\F_p$ is $\lceil \sqrt{p}\rceil$.
	\end{corollary}
	
	% In view of Proposition \ref{satbound}, we have the following. 
	
	% \begin{corollary}
	% Suppose that    
	% $M=2^{2n}+2^n+1$ and  $D'$ is a Singer  $(M,2^n+~1,1)$-difference set of the cyclic group $(\mathbb{Z}_M,+)$. 
	% Then $D'$ is $(2,-1)$-saturating set which attains the lower bound of Proposition \ref{satbound}.
	% \end{corollary}
	
	In general one can prove the following, along the same lines.
	
	\begin{prop}
		If $D$ is a $(v,k,\lambda)$-difference set in the group $G$ with numerical multiplier $2$ then $D$ saturates the $3$-$\mathrm{AP}$s.  Moreover, if $\lambda=1$ also holds, then $D$ is a complete $3$-$\mathrm{AP}$ free set.%\todo{Lehet hogy csak a Singer ilyen csak nem tudom, hogy ez bizonyitott-e.} %In particular if $q$ is a $2$-power then there exist complete $3$-$\mathrm{AP}$ free subsets of $\mathbb{Z}_{q^2+q+1}$ of small size.\todo{utolso mondat kuka}
	\end{prop}
	
	\begin{prop}
		If $D$ is a $(k^2+k+1,k+1,1)$-difference set in $G$, $0 \in D$, then $D$ is complete $(1,-1)$-avoiding of size $k+1$ and hence its size reaches the lower bound in Proposition \ref{satbound}. It follows that $a((1,-1),\mathbb{Z}_{k^2+k+1})=(k+1)$ if $k$ is a prime power.
	\end{prop}
	\begin{proof}
		By definition if $x\in G \setminus \{0\}$ then there exist $y,z \in D$ such that $x=y-z$. 
		
		Assume to the contrary $x=y-z$ for some pairwise distinct $x,y,z \in D$. Then $x-0=y-z$, contradicting the fact that $D$ is a $(k^2+k+1,k+1,1)$-difference set.
		
		The last part follows from the existence of Singer-difference sets. 
	\end{proof}
	
	Note that $0\in D$ can always be obtained since translates of $D$ are difference sets as well.
	
	
	\subsection{\texorpdfstring{$(1/2,1/2)$-saturating sets in cyclic groups via additive bases}{(1/2,1/2)-saturating sets in cyclic groups via additive bases}}
	
	%\todo[inline]{Megj.: A small-t neha definialjuk, neha meg csak mint jelzo hasznaljuk.}
	\begin{comment}
	In this subsection, we borrow saturation type results for integer sequences to improve our upper bounds on $\sat(3-\mathrm{AP}, \F_q)$ and $\sat(W, \F_q)$.
	
	
	We start with upper bounds on $\sat(W, q)$ for $W=\{(2,-1)\}$, which, as mentioned before, also serves as upper bound for $\sat(3-\mathrm{AP}, q)$.
	Fang proved that $S\cap [1, M]$ $(2, -1)$-saturates $[1,M]$ for all positive integer $M$.  More precisely, this is a direct consequence of his proof for $M\geq 32$ while it can be checked for all smaller positive integers. Fang also proved that  $\limsup_{M \to \infty}\frac{|S\cap [1, M]|}{\sqrt|M|}\leq \sqrt{12.8}$. This has a direct consequence for saturation in cyclic groups of prime order.
	
	
	\begin{prop}\label{integer} For every large enough prime $p$,  the size of the smallest $(2,-1)$-saturating set in $\F_p$ is at most $3.578\sqrt{p}$.
	\end{prop}
	
	\begin{proof}
	Take the integer sequence of Chen and Fang \cite{Chen, Fang}, given by $F=\bigcup _{l=0}^\infty T_l,$ where
	$$T_l=\{u4^l+v_{l-1}4^{l-1}+\cdots +v_04^0:
	% \hspace{0.5cm}  $
	v_i\in \{2,3\}\text{ for }i=0,1,\ldots ,l-1, u\in \{1,2,3,4\}\}.$$
	
	Let us intersect this with $[1,n]$ where $n=p$. This set $S$ consists of at most $3.578\sqrt{p}$ integers of $[1,n]$.
	We consider its elements as elements of $\F_p$, and want to show that for  each $x \in \F_p \setminus S$  there is a $3$-$\mathrm{AP}$ of $\F_p$ consisting of $x$ and two  elements of $S$. Since Chen and Fang proved that in $[1,p]$, there exist $2$ elements $a, a'$ of $S$ s.t. $a<a'<x$ and $a, a', x$ form a $3$-term arithmetic progression, we are done.
	\end{proof}
	
	In Theorem \ref{2-1sat}, we improved the constant $3.578$ from Proposition \ref{integer}, and proved that in fact the constant can be chosen to be smaller than $3$. 
	In order to do so, we describe the set of negative integers which are $(2, -1)$ saturated by  a KEZDOSZELET\todo{mi ez? (Z) Fu nem tudom, esetleg keplettel kiirhatjuk (B)} of
	$$S=\{k=k_l4^l+k_{l-1}4^{l-1}+\cdots +k_14^{1}+k_04^0 \mid   k_i\in \{2,3\} \ \forall \  i\leq l-1, k_l\in \{1,2,3,4\} , l\in \mathbb{N}\}.$$
	%\begin{lemma}\label{neg}
	%   Suppose that $$S=\{k=k_l4^l+k_{l-1}4^{l-1}+\cdots +k_14^{1}+k_04^0 \mid   k_i\in \{2,3\} \ \forall \  i\leq l-1, k_l\in \{1,2,3,4\} , l\in \mathbb{N}\}.$$
	%Then for every  $n\in [-M, 0] $, there exist  \todo{exists $\alpha$?} and $a, b\in S\cap [1, \alpha \cdot M]$ such that 
	%$n, a, b$ form a 3-term arithmetic progression.
	%\end{lemma}
	%The proof below follows the ideas of Fang \cite{Fang}, who 
	
	%\begin{proof}
	Let $n$ be given in the form 
	$$n= -c\cdot 4^t+\sum_{i=0}^{t-1}n_i\cdot 4^i $$ for some integer $t$, and $c, n_i\in \{1,2,3,4\}$. We will determine
	$a$ and $b$ in the form
	$$ a=a_t4^t+a_{t-1}4^{t-1}+\cdots +a_14^{1}+a_04^0, \ b=b_t4^t+b_{t-1}4^{t-1}+\cdots +b_14^{1}+b_04^0$$ such that $n-2a+b=0,$ the digits $a_i$ and $b_j$ of $a$ and $b$ are permitted and   $b$ is bounded from above as required.
	Let $a_t=0$ and $b_t=c$, furthermore let $a_i$, $b_i$ be determined from $n_i$ for all $i\leq t-1$ according to Table \ref{tab:my_label}. %\todo{nem értem miért ilyen hülyén számozza. Igy jo? B Köszi!Z}
	
	\begin{table}[h!]\label{aibi}
	\centering
	\begin{tabular}{c||c|c|c|c|}
	$n_i$   & 1 & 2 & 3 & 4\\ \hline
	$a_i$  & 2 & 2 & 3 & 3 \\ \hline
	$b_i$  & 3 & 2 & 3 & 2
	\end{tabular}
	\caption{The value of $a_i$ and $b_i, i\leq t-1$, determined by the value of $n_i$.}
	\label{tab:my_label}
	\end{table}   
	
	
	It is easy to see that the chosen digits are admissible and that $n, a, b$ form a $3$-term arithmetic progression, since their digits do so.
	Consequently, every negative integer $n$ greater than or equal to $-c\cdot 4^t$ is saturated by  $$S\cap [1, c\cdot 4^t+3\sum_{i=0}^{t-1} 4^i],$$
	for all $c\in \{1,2,3,4\}$.
	
	As a consequence, we in turn have
	
	\begin{lemma}\label{alul} \ \\
	If  $M\in [2\cdot 4^l, 3\cdot 4^l)$, then $S\cap [1,M]$ $(2,-1)$-saturates $[- 4^l,0]$.\\
	If  $M\in [3\cdot 4^l, 4\cdot 4^l)$, then $S\cap [1,M]$ $(2,-1)$-saturates $[-2\cdot 4^l,0]$.\\
	If  $M\in [4\cdot 4^l, 5\cdot 4^l)$, then $S\cap [1,M]$ $(2,-1)$-saturates $[-3\cdot 4^l,0]$.\\
	If  $M\in [5\cdot 4^l, 8\cdot 4^l)$, then $S\cap [1,M]$ $(2,-1)$-saturates $[-4\cdot 4^l,0]$.
	\end{lemma}
	
	
	%This in turn implies that for all $M\in \left[\lfloor (c+\frac{1}{3})4^t\rfloor, \lfloor (c+\frac{4}{3})4^t\rfloor\right]$, the size of the interval of integers saturated by $S\cap [1,M]$, taking into consideration the saturated value of $0$ as well,  is at least $$M+c\cdot4^t+1\geq M(1+\frac{c}{(c+\frac{4}{3})})\geq \frac{10}{7}M.$$.
	
	
	Now we describe the largest number $n(M)$ for which  $S \cap [1, M]$ $(2, -1)$-saturates the whole integer set $[1, n(M)]$ with
	the sequence $S$ of Fang.
	Using his observations, we have the following Table \ref{tab:digi} on the leading digits $a_l$, $b_l$ of the numbers $a(n)$ and $b(n)\in S$ for which $a(n), b(n), n$ form
	a $3$-$\mathrm{AP}$. $n$ is given as $2\cdot 4^l\leq n \leq 8\cdot 4^l$.
	
	\begin{table}[h!]\label{digits}
	\centering
	\begin{tabular}{c||c|c|c|c|c|c|c|}
	$\mu(n) \in $  & [2, $2\frac13$) & [$2\frac13$,  $3\frac13$) & [$3\frac13$,  $4\frac13$) & [$4\frac13$, $5\frac13$) & [$5\frac13$,  $6\frac13$) & [$6\frac13$, $7\frac13$) & [$7\frac13$, 8) \\ \hline
	$a_l$  & 1 & 0 &   1 & 0  & 1 & 0 & 1\\ \hline
	$b_l$ & 1 &  1 &  2 & 2 & 3& 3& 4 
	\end{tabular}
	\caption{The value of $a_l$ and $b_l$, depending on $\mu(n)$, where $\mu(n)$ is defined by $\frac{n}{4^l}$.}
	\label{tab:digi}
	\end{table}   
	
	This implies that  the following holds.
	\begin{lemma}\label{felul}
	If  $M\in [(u+1)4^l, (u+2)4^l)$, then $n(M)\geq (2u+\frac{4}{3})4^l$ for all $u\in \{1,2,3\}$.  
	\end{lemma}
	\begin{proof}
	Take an arbitrary value $u+1> 1$ from the row corresponding to $b_l$. It means that in each column with $b_l\leq u$, every $n$ in view can be saturated with two numbers from which the larger one (i.e., $b$) is at most $u\cdot4^l+3\cdot \sum_{j=0}^{l-1}4^j<(u+1)\cdot 4^l$.
	\end{proof}
	
	Now we are ready to prove the main result in this subsection. 
	
	\begin{prop}\label{osszesen} \ \\
	If  $M\in [2\cdot 4^l, 3\cdot 4^l)$, then $S\cap [1,M]$ $(2,-1)$-saturates $[- 4^l,3\frac13\cdot 4^l)$.\\
	If  $M\in [3\cdot 4^l, 4\cdot 4^l)$, then $S\cap [1,M]$ $(2,-1)$-saturates $[-2\cdot 4^l,5\frac13\cdot 4^l)$.\\
	If  $M\in [4\cdot 4^l, 5\cdot 4^l)$, then $S\cap [1,M]$ $(2,-1)$-saturates $[-3\cdot 4^l,7\frac13\cdot 4^l)$.\\
	If  $M\in [5\cdot 4^l, 8\cdot 4^l)$, then $S\cap [1,M]$ $(2,-1)$-saturates $[-4\cdot 4^l,8\cdot 4^l)$.
	\end{prop}
	\begin{proof}
	A direct  application of Lemma \ref{alul} and Lemma \ref{felul}.   
	\end{proof}
	\begin{cor}
	For every  integer $M$, the length of the $(2, -1)$-saturated interval by $S\cap [1,M]$ is at least $\frac{13}9M$.  
	\end{cor}
	
	
	
	
	\begin{theorem}\label{2-1sat}
	$\sat((2,-1), p)< 2.977\sqrt{p}$ if $p$ is large enough.
	\end{theorem}
	\begin{proof}[Proof of Theorem \ref{2-1sat}]
	Choose the smallest integer $M$ for which $\frac{13}9M>p$ and take $S\cap [1,M]\in \F_p$. Then Proposition \ref{osszesen} implies that $S$ saturates the whole set $\F_p$, and $S$ is of size at most $\sqrt{12.8M}<\sqrt{12.8\cdot \frac{9}{13}p}+\varepsilon=\sqrt{\frac{64\cdot9}{65}p}+\varepsilon\approx 2.977\sqrt{p}$ if  $p$ is large enough. Here $\varepsilon<1$ if $p\geq 7$, and $\varepsilon\to 0$ if $p\to \infty.$
	\end{proof}
	
	
	\begin{remark}
	Depending on the representation of $p$ in quaternary numeral system, Proposition \ref{osszesen} implies that if $M$ belongs to a suitable range, the length of the $(2, -1)$-saturated interval can be almost as long as $\frac{10\frac{1}{3}}{4}M$, hence for such an $\{M,p\}$ pair with $\frac{10\frac{1}{3}}{4}M>p$, we have  $\sat((2,-1), p)< 2.226\sqrt{p}$.
	\end{remark}
	
	
	%In order to bound $f(n)$ depending on $\mu(n)$, observe that it is enough to consider $\min \frac{n}{b(n)}$ for a fixed $l$, which can be bounded as follows.
	%$$\min \left\{\frac{n}{b(n)} \mid n=\mu(n)\cdot4^l\right\}=\min \{1+\frac{b(n)-a(n)}{b(n)}\}\geq 1+\frac{(\frac{2}{3}+b_l)-(a_l+1)}{\frac{2}{3}+b_l},$$ apart from the case $\mu(n)\in [2, 2\frac{1}3]$ when this minimum can be arbitrary close to $1$ as $l$ tends to infinity.
	%Indeed, $\frac{b(n)-a(n)}{b(n)}$ is minimal if $a$ is maximal while $b$ is minimal, but we must take into account the condition $b>a$. Here we applied the fact that every digit of $b$ is at least $2$ and every digit of $a$ is at most $3$, apart from the leading digit. Moreover, $$ 4^k+4^{k-1}+\ldots +4+1<\frac134^{k+1}.$$
	
	%\end{proof}
	
	%\begin{cor}
	%   Lemma  \ref{neg} in turn shows that a set of integers of size asymptotically equal to $\sqrt{15M}$ can saturate the cyclic group of order $\frac{10}{7}M$, hence asymptotically $\sqrt{10.5q}$ elements suffice in $\mathbb{Z}_q$.
	%\end{cor}
	
	\end{comment}
	
	We continue with upper bounds on $\sat(W,\F_p)$ for $W=\{(1/2, 1/2)\}$.
	
	Here we refer to a construction which provides a good upper bound for the solution of the postage stamp problem which is very closely related to finite additive basis, see \cite{Nath, Habsi, Hof}. Recall that the problem was described in the Introduction.
	%The famous postage stamp problem can be phrased as follows.
	%Given an integer $m$ and a set $A$ of positive integers, find the smallest integer  that cannot be written as the sum $\sum_{i=1}^k a_i$ for some number $k \leq m$ of not necessarily distinct elements of $A$.
	
	
	%\todo[inline]{We should look up the results concerning bounds for finite additive bases, which corresponds to $(0.5, 0.5)$-saturation. See \cite{Nath, Hof, Habsi}. Here the bounds are even more sharp!} \todo[inline]{If $B$ is an additive basis then $2B \cup \{0\}$ will saturate the 3-\mathrm{AP}s}
	
	Let $[a,(t), b]$ denote $\{a+t\cdot h: h\in \mathbb{Z}\}\cap [a,b]$.
	\begin{constr}[Mrose, \cite{Mrose}]
		\label{Mr}
		For an arbitrary positive integer $t$ take a set $S$ of $7t+2$ elements as $S=\bigcup_{j=1}^5 A^{(j)}$, where\\ 
		$A^{(1)}:=[0, (1), t],$\\ $ A^{(2)}:= [2t, (t), 3t^2+t],$\\ $ A^{(3)}:=[3t^2+2t, (t+1), 4t^2+2t-1],$\\ $ A^{(4)}:=[6t^2+4t, (1), 6t^2+5t],$\\ $A^{(5)}:=[10t^2+7t, (1), 10t^2+8t]$. 
	\end{constr}
	
	\begin{prop}[Mrose, \cite{Mrose}] \label{Mrose1}
		$(S+S)\supset [0, 14t^2+10t-1]$ holds for the Mrose construction $S$ with parameter $t$. 
	\end{prop}
	
	We apply this classical construction to prove the following upper bound for $\sat(3-\mathrm{AP}, \F_p)$ via $W$-saturation for $W=\{(1/2, 1/2)\}$.
	
	\begin{prop} Suppose that $m$ is odd. Then
		\[\sat(3-\mathrm{AP},\mathbb{Z}_m) \leq \sat((1/2, 1/2), \mathbb{Z}_m)\leq (\sqrt{3.5}+o(1))\sqrt{m}\approx 1.87\sqrt{m}.\]
	\end{prop}
	
	
	\begin{proof} Choose the least integer $t$ such that $14t^2+10t-1\geq m$ holds, i.e., $$14t^2+10t-1\geq m\ge 14(t-1)^2+10(t-1)-1.$$ Consider the set $S$  $\pmod m$ obtained in Construction \ref{Mr}. Since $S+S= \mathbb{Z}_m$, we also have $$\left\{\frac{s}{2}+\frac{s'}{2} \  : \  s, s'\in S\subset \mathbb{Z}_m\right\}=\mathbb{Z}_m,$$ since $\gcd(2,m)=1$. 
		Note that for $x\notin S$ and $x=s/2+s'/2$, $s,s' \in S$, we cannot have $s=s'$ and hence $x$ is saturated by two distinct elements of $S$. 
		Hence $S$ is a $(1/2, 1/2)$-saturating set in $\mathbb{Z}_m$ of size $|S|=7t+2$ while $m\ge 14t^2-18t+3> \frac{2}{7}|S|^2-4|S|.$ From this, we get that $\sat((1/2, 1/2), \mathbb{Z}_m)< 7+\sqrt{49+3.5m}.$ 
	\end{proof}


	
	\subsection{\texorpdfstring{Complete $(2,-1)$-avoiding and $(2,-1)$-saturating sets in cyclic groups}{Complete (2,-1)-avoiding and (2,-1)-saturating sets in cyclic groups}}
	
	We will say that $S$ $(2,-1)$-saturates $[x,y]$ if for each $z\in [x,y]$ there exist $a,b\in S$ such that $z=2a-b$. 
	
	\begin{remark}
		Every integer $1\leq k  \leq \frac43 (4^{n}-1)$ can be written in a unique way as 
		\[k=k_{l}4^{l}+\cdots +k_0 4^0,\]
		where $k_i\in \{1,2,3,4\}$ and $0 \leq l\leq n-1$. %Note that $k_i$ is the residue of $k-k_0-k_14-\ldots-k_{i-1}4^{i-1}$ modulo $4^i$ such that in case of residue equal to $0$ (and $k\neq k_0+k_14+\ldots+k_{i-1}4^{i-1}$) we put $k_i=4$ instead of $k_i=0$.
		This representation of positive integers is known as the bijective base-4 numeral system, see e.g, \cite{Smullyan}.
	\end{remark}
	
	\begin{comment}
	
	\todo[inline]{Legyen inkább bijective base-4 numeral system. \url{en.wikipedia.org/wiki/Bijective_numeration}.\\ Here comes a reshuffled version. (Z)}
	
	\bigskip
	
	\hline
	
	\bigskip
	
	\begin{constr}\label{zm}
	Let
	\[H_l=\{2\cdot4^l+v_{l-1}4^{l-1}+\cdots +v_04^0 : 
	v_i\in \{2,3\}\text{ for }i=0,1,\ldots ,l-1\},\]
	and for a non-negative integer $r< l$, let
	\[L_{r,l}=\{3\cdot4^l+v_{l-1}4^{l-1}+\cdots +v_04^0 : 
	v_i\in \{2,3\}\text{ for }i\le l-1-r\}, v_i=2\text{ for }i\in [l-r, l-1]\}.\]
	
	Let
	\[K_l=\{2\cdot4^l+v_{l-1}4^{l-1}+\cdots +v_04^0 : 
	v_i\in \{1,2,3,4\}\text{ for }i=0,1,\ldots ,l-1\}.\]
	For a positive integer $r\le l$, let
	\[R_{r,l}=\{4^l+v_{l-1}4^{l-1}+\cdots +v_04^0 : 
	v_i\in \{1,2,3,4\}\text{ for }i\le l-1-r\}, v_i=4\text{ for }i\in [l-r, l-1]\}.\] Finally 
	\[S_l=\{v_l\cdot4^l+v_{l-1}4^{l-1}+\cdots +v_04^0 : 
	v_i\in \{1,2,3,4\}\text{ for }i=0,1,\ldots ,l\},\] so $S_l$ is the set of integers with exactly $l+1$ digits in the bijective base $4$ number system. 
	\end{constr}
	
	\begin{prop}\label{halmazok}
	\begin{enumerate}[\rm(1)]
	\item $H_l$ $(2,-1)$-saturates any subset of $K_l$. 
	\item $H_l\cup L_{r,l}$ $(2,-1)$-saturates any subset of $K_l\cup R_{r,l}$ for every positive integer $r<l$. 
	\item  $H_l\cup L_{0,l}$ $(2,-1)$-saturates any subset of $S_l$ for every positive integer $r<l$. 
	\end{enumerate}
	\end{prop}
	
	
	\begin{proof}
	
	Let $k$ have $l+1$ digits, then consider the $l$-digit numbers $a$ and $b$ from the saturating set defined above, where the $i$th digit of $a$ and $b$ is defined according to Table \ref{tab:my_label1}. Observe that when $k_i=4$, we have $b_i=2$, this will ensure that the elements of $R_{r,l}$ are indeed saturated by a pair $(a,b)\in H_l \times L_{r,l}$.
	Note that $k_i=2a_i-b_i$ and hence 
	\[k=(2a_{t-1}-b_{t-1})4^{t-1}+\cdots +(2a_0-b_0) 4^0=\]
	\[2 \sum_{i=0}^{t-1} a_{i}4^i - \sum_{i=0}^{t-1} b_i 4^i=2a-b,\]
	with $a,b \in S$. It follows that $b, a, k=2a-b$ is a $3$-$\mathrm{AP}$ with difference $a-b$. 
	\end{proof}
	
	\begin{theorem}\label{4pergyok5}
	Let $m$ be of form $m=\alpha\cdot 4^l$ for $\alpha\in (1, 4]$.
	\end{theorem}
	
	\begin{proof}
	Observe that $|H_l|=2^l$ and $|L_{r,l}|=2^{l-r}$, while $|K_l|=4^l$, $|R_{r,l}|=4^{l-r}$, $|S_l|=4^{l+1}$ holds.%\todo[inline]{to be continued, but hopefully clear}
	
	If $\alpha\in (1+\frac{1}{4^{r+1}}, 1+\frac{1}{4^r}]$ for a positive integer $r$, then take an interval $I$ of $m$ elements in $K_l\cup R_{r,l}$, which will be $(2,-1)$ saturated by $H_l\cup L_{r,l}$, hence the size of the saturating set is $2^l+2^{l-r}=\frac{(1+2^{-r})}{\sqrt{\alpha}}\sqrt{m}.$ Since $\min{\alpha}=1+\frac{1}{4^{r+1}}$ in this case, we have an upper bound $\frac{1+2^{-r}}{\sqrt{1+{4^{-r-1}}}}\sqrt{m}$ for the saturating set.\\
	If $\alpha\in ( 1+\frac{1}{4}, 4]$ then take an interval $I$ of $m$ elements in $S_l$, which will be $(2,-1)$ saturated by $H_l\cup L_{0,l}$, hence the size of the saturating set is $2^l+2^{l}=\frac{2}{\sqrt{\alpha}}\sqrt{m}.$ Since $\min{\alpha}=1+\frac{1}{4}$ in this case, we have an upper bound $\frac{4}{\sqrt{5}}\sqrt{m}$ for the saturating set.    
	\end{proof}
	
	\hline
	
	\begin{theorem}
	\label{gyok3}
	Put 
	\[H_l=\{v_{l-1}4^{l-1}+\cdots +v_04^0 : 
	v_i\in \{2,3\}\text{ for }i=0,1,\ldots ,l-1\},\]
	and put
	\[K_l=\{v_{l-1}4^{l-1}+\cdots +v_04^0 : 
	v_i\in \{1,2,3,4\}\text{ for }i=0,1,\ldots ,l-1\},\]
	so $K_l$ is the set of integers with exactly $l$ digits in the bijective base $4$ number system. 
	
	Note that $K_l=[\frac13(4^l-1),\frac43(4^l-1)]$, so $|K_l|=4^l$.
	
	The smallest integer of $H_l$ is $\frac23(4^l-1)$, the largest one is $4^l-1$ and $|H_l|=2^l$.
	
	\begin{enumerate}[\rm(1)]
	\item for $1 \leq i\leq j$, $H_i \cup H_{i+1} \cup \ldots \cup H_{j}$ $(2,-1)$-saturates any subset of $K_i \cup K_{i+1} \cup \ldots \cup K_{j}$.
	\end{enumerate}
	
	
	% Then $S$ is complete $3$-$\mathrm{AP}$ free in $I=[1,\frac43 (4^n-1)]$ of size $2^{n+1}-2$.
	
	For some positive integer $n$, let $m$ denote an integer such that $4^{n-1} < m \leq 4^n$.
	
	\begin{enumerate}[\rm(2)]
	\item If $m=4^n$, then consider the elements of $H_n$ and $K_n$ as representatives for the elements of $\mathbb{Z}_{4^n}$. The $2^n$ elements corresponding to $H_n$ form a complete $(2,-1)$-avoiding set in $\mathbb{Z}_m$.
	
	
	
	\item[\rm(3)] If $\frac13(4^n-1)+1\leq m<4^n$ then consider any subset $H_n \subseteq [x,y] \subseteq K_n$ of size $m$ as a representative for $\mathbb{Z}_m$. Then the elements corresponding to $H_n$ form a $(2,-1)$-saturating set of size at most $\sqrt{3 m}$ in $\mathbb{Z}_m$.
	
	If $\frac23(4^n-1) \leq m$ then $H_n$ corresponds to a complete $(2,-1)$-avoiding set.
	
	\item[\rm(4)] If $4^{n-1}<m\leq \frac13(4^n-1)$, then choose $[0,m]$ as a representative for $\mathbb{Z}_m$. Then $\{0\}\cup H_1 \cup H_2 \cup \ldots \cup H_{n-1} \subseteq [0,m]$ and it corresponds to a $(2,-1)$-saturating set of $\mathbb{Z}_m$ of size $2^{n}-1 \leq 2\sqrt{m}$. \todo[inline]{By choosing a better subset (alkalmas fels\H o kezd\H o szeleteket) I think it is possible to improve the constant to $\sqrt{3}$. Az biztos, hogy a $2$-vel mint konstanssal m\H ukodik. Teh\'at jobbat fog adni, mint Remark 4.4-ig az el\H oz\H o fejezet eredmenyei.}
	\end{enumerate}
	
	% Note that
	% \[|S| \approx \sqrt{3}\sqrt{|I|}\]
	% and hence for an integer $s\leq \frac43(4^n-1)$ close to $\frac43(4^n-1)$ we obtain a $(2,-1)$-saturating set in $\mathbb{Z}_s$ of size roughly $\sqrt{3}\sqrt{s}$ 
	% (after adding $0$ to $S \cap [1,s]$ if necessary).
	\end{theorem}
	\begin{proof}
	First we show that $S$ is $3$-$\mathrm{AP}$-free.
	Suppose to the contrary that $a<b<c$ are three elements of $S$ forming an arithmetic progression. Assume that $r$ is maximal such that the coefficients $a_r,b_r,c_r$ of $4^r$ are not all equal in the expressions of $a$, $b$, $c$ as above. Then clearly $a_r\leq b_r \leq c_r$. 
	If $a_r=b_r=2$, then $c_r=3$ and $b-a$ is at most $4^{r-1}+\ldots+1$, while $c-b$ is at least $4^r-4^{r-1}-\ldots-1$. It follows that $b-a \neq c-b$. Similar arguments work when $a_r=2$, $b_r=c_r=3$ and when $a_r=0$.
	
	Next we show that $S$ is $(2, -1)$-saturating.
	Consider $k\in [1,\frac43(4^n-1)]$, then
	\[k=k_{l}4^{l}+\cdots +k_0 4^0 : k_i\in \{1,2,3,4\}\text{ for }i=0,1,\ldots ,l,\]
	for some $l\leq n-1$. If $k$ has $t$-digits (so $t$ is maximal such that $k_{t-1}$ is non-zero), then consider the $t$-digit numbers $a$ and $b$ according to the following table:
	\begin{table}[h!]
	\centering
	\begin{tabular}{c||c|c|c|c|}
	$k_i$   & 1 & 2 & 3 & 4\\ \hline
	$a_i$  & 2 & 2 & 3 & 3 \\ \hline
	$b_i$  & 3 & 2 & 3 & 2
	\end{tabular}
	\caption{The value of $a_i$ and $b_i, i\leq t-1$, determined by the value of $k_i$.}
	\label{tab:my_label1}
	\end{table}
	Note that $k_i=2a_i-b_i$ and hence 
	\[k=(2a_{t-1}-b_{t-1})4^{t-1}+\cdots +(2a_0-b_0) 4^0=\]
	\[2 \sum_{i=0}^{t-1} a_{i}4^i - \sum_{i=0}^{t-1} b_i 4^i=2a-b,\]
	with $a,b \in S$. It follows that $b, a, k=2a-b$ is a $3$-$\mathrm{AP}$ with difference $a-b$. 
	\end{proof}
	\todo[inline]{
	a((2,-1),G)*a((2,-1),H) >= a(3-AP,G×H)
	
	(And this motivates further the study of a((2,-1),G), even thought it is interestong on its owm right since "van hogy eles az also becsles".)}
	% \begin{theorem}
	% \label{gyok3}
	%     Put $S=\bigcup _{l=1}^n H_l$, where
	%  \[H_l=\{v_{l-1}4^{l-1}+\cdots +v_04^0 : 
	% v_i\in \{2,3\}\text{ for }i=0,1,\ldots ,l-1\}.\]
	% Then $S$ is complete $3$-$\mathrm{AP}$ free in $I=[1,\frac43 (4^n-1)]$ of size $2^{n+1}-2$.
	
	% Note that
	% \[|S| \approx \sqrt{3}\sqrt{|I|}\]
	% and hence for an integer $s\leq \frac43(4^n-1)$ close to $\frac43(4^n-1)$ we obtain a $(2,-1)$-saturating set in $\mathbb{Z}_s$ of size roughly $\sqrt{3}\sqrt{s}$ 
	% (after adding $0$ to $S \cap [1,s]$ if necessary).
	% \end{theorem}
	% \begin{proof}
	%     First we show that $S$ is $3$-$\mathrm{AP}$-free.
	%     Suppose to the contrary that $a<b<c$ are three elements of $S$ forming an arithmetic progression. Assume that $r$ is maximal such that the coefficients $a_r,b_r,c_r$ of $4^r$ are not all equal in the expressions of $a$, $b$, $c$ as above. Then clearly $a_r\leq b_r \leq c_r$. 
	%     If $a_r=b_r=2$, then $c_r=3$ and $b-a$ is at most $4^{r-1}+\ldots+1$, while $c-b$ is at least $4^r-4^{r-1}-\ldots-1$. It follows that $b-a \neq c-b$. Similar arguments work when $a_r=2$, $b_r=c_r=3$ and when $a_r=0$.
	
	% Next we show that $S$ is $(2, -1)$-saturating.
	%  Consider $k\in [1,\frac43(4^n-1)]$, then
	%     \[k=k_{l}4^{l}+\cdots +k_0 4^0 : k_i\in \{1,2,3,4\}\text{ for }i=0,1,\ldots ,l,\]
	% for some $l\leq n-1$. If $k$ has $t$-digits (so $t$ is maximal such that $k_{t-1}$ is non-zero), then consider the $t$-digit numbers $a$ and $b$ according to the following table:
	%      \begin{table}[h!]
	%      \centering
	%      \begin{tabular}{c||c|c|c|c|}
	%         $k_i$   & 1 & 2 & 3 & 4\\ \hline
	%         $a_i$  & 2 & 2 & 3 & 3 \\ \hline
	%         $b_i$  & 3 & 2 & 3 & 2
	%      \end{tabular}
	%      \caption{The value of $a_i$ and $b_i, i\leq t-1$, determined by the value of $k_i$.}
	%      \label{tab:my_label1}
	%  \end{table}
	%  Note that $k_i=2a_i-b_i$ and hence 
	%  \[k=(2a_{t-1}-b_{t-1})4^{t-1}+\cdots +(2a_0-b_0) 4^0=\]
	%  \[2 \sum_{i=0}^{t-1} a_{i}4^i - \sum_{i=0}^{t-1} b_i 4^i=2a-b,\]
	%  with $a,b \in S$. It follows that $b, a, k=2a-b$ is a $3$-$\mathrm{AP}$ with difference $a-b$. 
	% \end{proof}
	
	\end{comment}
	
	
	
	
	\begin{constr}\label{z_m}
		
		Let
		\[H_l=\{v_{l-1}4^{l-1}+\cdots +v_04^0 : 
		v_i\in \{2,3\}\text{ for }i=0,1,\ldots ,l-1\},\]
		and
		\[K_l=\{v_{l-1}4^{l-1}+\cdots +v_04^0 : 
		v_i\in \{1,2,3,4\}\text{ for }i=0,1,\ldots ,l-1\},\]
		so $K_l$ is the set of integers with exactly $l$ digits in the bijective base-$4$ numeral system. 
	\end{constr}
	
	Note that $K_l=[\frac13(4^l-1),\frac43(4^l-1)]$, so $|K_l|=4^l$.
	
	The smallest integer of $H_l$ is $\frac23(4^l-1)$, the largest one is $4^l-1$, and $|H_l|=2^l$.
	
	\begin{theorem}
		\label{gyok3} \ \\ \vspace{-0.8cm}
		\begin{enumerate}[\rm(1)]
			\item The set $H_i \cup H_{i+1} \cup \ldots \cup H_{j}$ $(2,-1)$-saturates any subset of $K_i \cup K_{i+1} \cup \ldots \cup K_{j}$ for every pair of positive integers $i\leq j$.
		\end{enumerate}
		
		
		% Then $S$ is complete $3$-$\mathrm{AP}$ free in $I=[1,\frac43 (4^n-1)]$ of size $2^{n+1}-2$.
		
		Given a positive integer $n$, let $m$ denote an integer such that $4^{n-1} < m \leq 4^n$.
		
		\begin{enumerate}[\rm(2)]
			\item If $m=4^n$, then consider the elements of $H_n$ and $K_n$ as representatives for the elements of $\mathbb{Z}_{4^n}$. The $2^n$ elements corresponding to $H_n$ form a complete $(2,-1)$-avoiding set in $\mathbb{Z}_m$.
			
			
			
			\item[\rm(3)] If $\frac13(4^n-1)+1\leq m<4^n$ then consider any interval $[x,y]$,  $H_n \subseteq [x,y] \subseteq K_n$, of size $m$ as a representative for $\mathbb{Z}_m$. Then the elements corresponding to $H_n$ form a $(2,-1)$-saturating set of size less than $\sqrt{3 m}$ in $\mathbb{Z}_m$.
			
			If $\frac23(4^n-1) < m$ then $H_n$ corresponds to a complete $(2,-1)$-avoiding set.
			
			\item[\rm(4)] If $4^{n-1}<m\leq \frac13(4^n-1)$, then let $1\leq k\leq n-1$ be maximal such that 
			\[4^{n-1}< m \leq (4^n-4^{k-1})/3.\]
			Then $S:=H_{k-1} \cup H_{k}\cup \ldots \cup H_{n-1}$ $(2,-1)$-saturates $I:=[\frac13(4^{k-1}-1),\frac13(4^{k-1}-1)+m-1]$ and has size less than $\sqrt{3m}$. Considering $I$ as representatives for $\mathbb{Z}_m$ the same holds for the elements corresponding to $S$. 
		\end{enumerate}
		
	\end{theorem}
	\begin{proof}
		We start with proving $(1)$. It is enough to prove that $H_t$ saturates $K_t$. 
		If $k$ has $t$ digits, then consider the $t$-digit numbers $a$ and $b$ according to Table \ref{tab:my_label1}.
		
		\begin{table}[h!!]
			\centering
			\begin{tabular}{c||c|c|c|c|}
				$k_i$   & 1 & 2 & 3 & 4\\ \hline
				$a_i$  & 2 & 2 & 3 & 3 \\ \hline
				$b_i$  & 3 & 2 & 3 & 2
			\end{tabular}
			\caption{The value of $a_i$ and $b_i, i\leq t-1$, determined by the value of $k_i$.}
			\label{tab:my_label1}
		\end{table}
		
		Note that $k_i=2a_i-b_i$ and hence 
		\[k=(2a_{t-1}-b_{t-1})4^{t-1}+\cdots +(2a_0-b_0) 4^0=\]
		\[2 \sum_{i=0}^{t-1} a_{i}4^i - \sum_{i=0}^{t-1} b_i 4^i=2a-b,\]
		with $a,b \in H_t$. It follows that $b, a,$ and $ k=2a-b$ form a $3$-$\mathrm{AP}$ with difference $a-b$. 
		
		To prove $(2)$ and $(3)$ first we show that $S:=\cup_{i=1}^{\infty} H_i$ is $3$-$\mathrm{AP}$ free in $\mathbb{Z}$.
		Suppose to the contrary that $a<b<c$ are three elements of $S$ forming an arithmetic progression. Assume that $r$ is maximal such that the coefficients $a_r,b_r,c_r$ of $4^r$ are not all equal in the expressions of $a$, $b$, $c$ as above. Then clearly $a_r\leq b_r \leq c_r$. 
		If $a_r=b_r=2$, then $c_r=3$ and $b-a$ is at most $4^{r-1}+\ldots+1$, while $c-b$ is at least $4^r-4^{r-1}-\ldots-1$. It follows that $b-a \neq c-b$. Similar arguments work when $a_r=2$, $b_r=c_r=3$ and when $a_r=0$.
		
		If we consider the set $K_n$ with modulo $4^n$ addition, then  the $(2,-1)$-saturation property clearly holds. We show that the  $(2,-1)$-avoiding property  holds as well. Suppose to the contrary that $a<b<c$ are three elements of $H_n$ forming an arithmetic progression when considered modulo $4^n$.
		Then for some difference $0<d<4^n$ we have $a\equiv c+d \pmod {4^n}$ and 
		either $c-b=d$, or $b-a=d$. 
		From $a\equiv c+d \pmod {4^n}$ it follows that 
		$d \geq 4^n+\frac23(4^n-1)-(4^n-1)=\frac23(4^n-1)+1$ and hence $d=c-b$ and $d=b-a$ are impossible because of $(4^n-1)-\frac23(4^n-1)=\frac13(4^n-1) \geq \max\{c-b,b-a\}$. This proves $(2)$.
		
		The first part of $(3)$ follows from the fact that 
		\[\sqrt{3m}\geq \sqrt{4^n+2}>2^n=|H_n|.\]
		
		To prove the second part suppose to the contrary that $a<b<c$ are three elements of $H_n$ forming an arithmetic progression when considered modulo $4^n$.
		Then for some difference $0<d<m$ we have $a\equiv c+d \pmod m$ and 
		either $c-b=d$, or $b-a=d$. 
		From $a\equiv c+d \pmod m$ it follows that 
		$d \geq m+\frac23(4^n-1)-(4^n-1)=m-\frac13(4^n-1)> \frac13(4^n-1) $ and hence $d=c-b$ and $d=b-a$ are impossible because of $(4^n-1)-\frac23(4^n-1)=\frac13(4^n-1) \geq \max\{c-b,b-a\}$.
		
		To prove $(4)$ assume that $1 \leq k \leq n-1$ is maximal such that 
		$3m \leq 4^n-4^{k-1}$. It follows that 
		\[3m > 4^n-4^k.\]
		First note that $S:=H_{k-1}\cup\ldots \cup H_{n-1}$ has size $2^{k-1}+\ldots+2^{n-1}=2^{k-1}(1+\ldots+2^{n-k})=2^{k-1}(2^{n-k+1}-1)=2^n-2^{k-1}$. Since $S$ $(2,-1)$-saturates $Z:=K_{k-1}\cup \ldots \cup K_{n-1}$ and 
		$I:=[\frac13(4^{k-1}-1),\frac13(4^{k-1}-1)+m-1]\subseteq Z$, it follows that $S$ $(2,-1)$-saturates $I$. We want to show $|S|< \sqrt{3m}$, that is,
		\[2^n-2^{k-1}<\sqrt{3m}.\]
		Clearly, it is enough to prove $4^n+4^{k-1}-2^{n+k}< 3m$, which follows from $4^n+4^{k-1}-2^{n+k}< 4^n-4^k < 3m$. 
	\end{proof}
	
	\subsection{Constructions in abelian groups of composite order}
	\label{last}
	
	% \todo[inline]{In $T$ it should also hold that $2x=2y$ does not have a non-trivial solution. Otherwise it is not true. ell: hol használjuk. maradjon ptn rend, vagy csak a Z2 kizárando}
	
	
	
	
	\begin{theorem}
		Let $G$ denote a commutative group, $H$ a subgroup of $G$. Put $S=\{a_1,a_2,\ldots,a_s\} \subseteq H$ of size $s$ and $T=\{b_1+H,b_2+H,\ldots,b_t+H\} \subseteq G/H$ of size $t$. Let $w_1, w_2 \in \mathbb{Z}$ such that
		$w_1+w_2=1$ holds. 
		Define
		\[X=\{a_i + b_j : i\in \{1,\ldots,s\},\, j \in\{1,\ldots,t\}\} \subseteq G.\]
		\begin{enumerate}[\rm(1)]
			\item If $S$ and $T$ are $(w_1,w_2)$-saturating in the groups $H$ and $G/H$, respectively, then
			$X$ is $(w_1,w_2)$-saturating in $G$.
			\item Assume that $S$ and $T$ are $(w_1,w_2)$-avoiding in the groups $H$ and $G/H$, respectively, and the order of $x-y$ is not a divisor of $w_1$ in the group $H$ ($G/H$) for each $x,y\in S$ (for each $x,y \in T$). Then $X$ is $(w_1,w_2)$-avoiding in $G$.
			\item Assume that $S$ and $T$ are complete $(w_1,w_2)$-avoiding in the groups $H$ and $G/H$, respectively, and the order of $x-y$ is not a divisor of $w_1$ in the group $H$ ($G/H$) for each $x,y\in S$ (for each $x,y \in T$). Then $X$ is complete $(w_1,w_2)$-avoiding in $G$.
		\end{enumerate}
	\end{theorem}
	
	\begin{proof} First suppose that $S$ and $T$ are $(w_1,w_2)$-saturating sets.
		Take some $c\in G$. Then $c=a+b$ where $a\in H$ and $b+H$ is an element of $G/H$.
		
		By the saturation property, there exist distinct $b_i+H, b_j+H \in T$ such that $w_1( b_i+H) + w_2 (b_j+H) =b+H$. 
		If $w_1b_i + w_2b_j = a'+ b$, for some $a'\in H$, then take some distinct  
		$a_f, a_g \in S$ such that $w_1a_f + w_2a_g=a- a'$.
		Then 
		\[w_1(a_f+b_i)+w_2(a_g+b_j)=c.\]
		If $a_f+b_i=a_g+b_j$, then $a_f-a_g=b_j-b_i \in H$, a contradiction since $i\neq j$. It follows that $X$ saturates $c\in G$. 
		
		Now assume that the conditions of part $(2)$ hold, and 
		\[w_1(a_f+ b_i)+w_2(a_g+b_j)= a_h+b_k\]
		for some elements $a_f+b_i$, $a_g+b_j$, $a_h+b_k$ of $X$. 
		Thus $w_1(b_i+H)+w_2(b_j+H)=b_k+H$ and hence $\{b_i+H, b_j+H, b_k+H\}$ is a set of size at most $2$. 
		If it has size $2$, then we may assume $b_i \neq b_k$. Then the order of $(b_i+H)-(b_k+H)$ is divisible by $w_1$, a contradiction. If $b_i+H=b_j+H=b_k+H$, then $w_1 a_f + w_2 a_g = a_h$. It follows that the size of $\{a_f,a_g,a_h\}$ is at most $2$. If it has size $2$, then we may assume $a_f \neq a_h$. Then the order of $a_f-a_h$ divides $w_1$, a contradiction. Consequently, $a_f=a_g=a_h$ holds and hence
		$a_f+b_i=a_g+b_j=a_h+b_k$. 
		
		The third part is a direct consequence of the first two.
	\end{proof}
	
	In the next result $\mathbb{Z}_r$ is considered as $\{0,1,\ldots,r-1\}$ with operation the usual addition in $\mathbb{Z}$ modulo $r$.
	
	\begin{corollary}
		\label{prodc} 
		Put $S=\{a_1,a_2,\ldots,a_s\} \subseteq \mathbb{Z}_m$ and $T=\{b_1,b_2,\ldots,b_t\} \subseteq \mathbb{Z}_n$. Let $w_1, w_2 \in \mathbb{Z}$ such that
		$w_1+w_2=1$ hold.
		Define
		\[X=\{a_i n + b_j : i\in \{1,\ldots,s\},\, j \in\{1,\ldots,t\}\} \subseteq \mathbb{Z}_{nm}.\]
		\begin{enumerate}[\rm(1)]
			\item If $S$ and $T$ are $(w_1,w_2)$-saturating, then
			$X$ is $(w_1,w_2)$-saturating in $\mathbb{Z}_{nm}$.
			\item Assume that $S$ and $T$ are $(w_1,w_2)$-avoiding, the difference of distinct elements of $S$ is not divisible by $m$, the difference of distinct elements of $T$ is not divisible by $n$. Then 
			$X$ is $(w_1,w_2)$-avoiding in $\mathbb{Z}_{nm}$.
			\item Assume that $S$ and $T$ are complete $(w_1,w_2)$-avoiding, the difference of distinct elements of $S$ is not divisible by $m$, the difference of distinct elements of $T$ is not divisible by $n$. Then
			$X$ is complete $(w_1,w_2)$-avoiding in $\mathbb{Z}_{nm}$.
		\end{enumerate}
	\end{corollary}
	
	\begin{example}
		$\{0,1,2\}$ is complete $(3,-2)$-avoiding in $\mathbb{Z}_9$ and hence  $a((3,-2),\mathbb{Z}_{9^n})=~3^n$.
	\end{example}
	
	\begin{example}
		$\{0,1\}$ is complete $(2,-1)$-avoiding in $\mathbb{Z}_4$ and hence  $a((2,-1),\mathbb{Z}_{4^n})=2^n$, as we already saw in the previous section.
	\end{example}
	
	\section{Concluding remarks and open problems}
	
	In this paper, we proved that in a large family of vector spaces, the minimum size of a complete $3$-$\mathrm{AP}$ free set is equal to a small absolute constant multiple of the  lower bound. However, it remained an open question to decide whether this is true for every vector space $\F_q^n$, $q>2$.
	
	\begin{problem}[Minimum size complete $3$-AP free sets] \ \\ Is it true that $$a(3-AP, \F_q^n)<C\cdot \sqrt{q^n}$$ for an absolute constant $C$, that is, the natural lower bound is tight up to a constant factor?
	\end{problem} 
	Concerning cyclic groups, we pose the following
	\begin{problem} Is it true that
		$a(3-AP,\mathbb{Z}_m)< a((2,-1), \mathbb{Z}_m)<\sqrt{c_m \cdot m}$ holds for $c_m\in [1, 1.5]$, a constant depending only on $m$, for all large enough values of $m$? 
	\end{problem}
	The first inequality follows from the definition (see Remark \ref{rem}), while we proved the second inequality for a dense set of natural numbers $m$ in Theorem \ref{maingroups}. It would be also interesting to see an improvement on the constant $c_m$.
	
	\section*{Acknowledgement}
	The authors would like to thank the referees for their helpful suggestions.\\
	This work was supported by the Italian National Group for Algebraic and Geometric Structures and their Applications (GNSAGA--INdAM). Both authors acknowledge the partial support of the Hungarian Research Grant (NKFI) K 124950. The first author is supported by the J\'anos Bolyai Research Scholarship of the Hungarian Academy of Sciences. The second author is supported by the Hungarian Research Grant (NKFI) No. PD  134953 and by the University Excellence Fund of Eötvös Loránd University.
	
	
	% \subsection{vanilla3}
	
	% In the next result $\mathbb{Z}_r$ is considered as $\{0,1,\ldots,r-1\}$ with operation the usual addition in $\mathbb{Z}$ modulo $r$.
	
	% \begin{theorem}
	% \label{prodc} 
	% Put $S=\{a_1,a_2,\ldots,a_s\} \subseteq \mathbb{Z}_m$ and $T=\{b_1,b_2,\ldots,b_t\} \subseteq \mathbb{Z}_n$. Let $w_1, w_2 \in \mathbb{Z}$ such that
	%    $w_1+w_2=1$ holds. 
	% Define
	% \[X=\{a_i n + b_j : i\in \{1,\ldots,s\},\, j \in\{1,\ldots,t\}\} \subseteq \mathbb{Z}_{nm}.\]
	% \begin{enumerate}[\rm(1)]
	% 	\item If $S$ and $T$ are $(w_1,w_2)$-saturating, then
	% 	$X$ is $(w_1,w_2)$-saturating in $\mathbb{Z}_{nm}$.
	% 	\item If $S$ and $T$ are $(w_1,w_2)$-avoiding, then
	% 	$X$ is $(w_1,w_2)$-avoiding in $\mathbb{Z}_{nm}$.
	% 	\item If $S$ and $T$ are complete $(w_1,w_2)$-avoiding, then
	% 	$X$ is complete $(w_1,w_2)$-avoiding in $\mathbb{Z}_{nm}$.
	% \end{enumerate}
	% \end{theorem}
	
	% \begin{proof} First suppose that $S$ and $T$ are $(w_1,w_2)$-saturating sets.
	% Take some $c\in \{0,1,\ldots,nm-~1\}$. Then $c=an+b$ for some uniquely determined $b\in \{0,1,\ldots,n-1\}$ and $a\in \{0,1,\ldots,m-~1\}$.
	
	% By the saturation property, there exist $b_i,b_j \in T$ such that $w_1b_i+w_2b_j=b +kn$ for some integer $k$. Moreover, there exists 
	% $a_f, a_g \in S$ such that $w_1 a_f+w_2 a_g \equiv a-k \pmod m$. It follows that 
	% \[w_1(a_f n + b_i)+w_2(a_g n + b_j)=(w_1a_f+w_2a_g)n+b+kn \equiv an+b \pmod {mn},\]
	% that is, $X$ saturates $c\in \mathbb{Z}_{nm}$. 
	
	% Now assume that $S$ and $T$ are $(w_1,w_2)$-avoiding and 
	% \[w_1(a_fn+b_i)+w_2(a_gn+b_j)\equiv a_hn+b_k \pmod {nm}.\]
	% It follows that $w_1b_i+w_2b_j\equiv b_k \pmod n$ and hence $b_i=b_j=b_k$ follows from the  $(w_1,w_2)$-avoiding property of $T$. Then $n(w_1 a_f + w_2 a_g - a_h) \equiv 0 \pmod {nm}$ and hence $w_1 a_f + w_2 a_g \equiv a_h \pmod m$. It follows that $a_f=a_g=a_h$, which proves the second part.
	
	% The third part is a direct consequence of the first two.
	% \end{proof}
	
	
	
	
	%\section{Concluding remarks and open problems}
	%\begin{itemize}
	
	%\item What can we say about $3$-$\mathrm{AP}$ saturation in groups of even order? The elementary abelian $\F_2^n$ case is not interesting, but already the cyclic case of even order seems to be non-trivial. For example what is the lower bound in this case? Or what about the $(1/2,1/2)$-saturation? (which should have a new name when the group order is even since $x+x=g$ might have more than one or zero solution in $x$)
	
	%\item Maybe we can find small examples for $(1,-1)$-avoiding/saturating sets among   $(k^2,k,0,k-1)$-almost difference sets.
	
	%\item Is there an absolute constant $c$ such that in each abelian group of order $n$ there is a $3$-AP-saturating set of size at most $c\sqrt{n}$.
	
	% \item a(2, -1), Fp azokra a p-kre amiket nem kezeltünk?? 
	
	% \item what if $-2$ is not a square element? $\F_q$-s where the lower bound can be strengthened.
	% \item k-AP
	% \item further results for $W$-avoiding and satu sets, w.r.t. the structure of $W$.
	% \item dim= odd case: $\F_q^3$-ben keresni kis pèldàkat $\beta$-ratio-free halmazra (ami m\H ukodik nem n\'egyzet $q$-ra is). P\'eld\'aul masodrend\H u felulet? (azt gondoljuk az\'ert, hogy nem ez a j\'o nagysagrend, hanem a randomized algo által adott ponthalmaz ami kerüli a $\beta$ ratiot, lényegében uyanazt adja, mint amit a véletlen konstrukcióban számoltunk.)
	%\end{itemize}
	%\section{\texorpdfstring{$x^q$}{xq}}
	
	
	{\footnotesize
		\begin{thebibliography}{pippo}
			
			
			%\bibitem{LuMaPoTr2014}
			%{\sc G. Lunardon, G. Marino, O. Polverino and R. Trombetti:} Maximum scattered linear sets of pseudoregulus type
			%and the Segre Variety ${\cal S}_{n,n}$, J.\ Algebr.\ Comb.\ \textbf{39} (2014), 807--831.
			
			\bibitem{Alon} Alon, N.,  Shapira, A. (2005). Linear equations, arithmetic progressions and hypergraph property testing. Theory of Computing, 1(1), 177--216.
			
			\bibitem{Anbar} Anbar, N., Bartoli, D., Giulietti, M., Platoni, I. (2014). Small complete caps from singular cubics. Journal of Combinatorial Designs, 22(10), 409--424.
			
			\bibitem{Bartoli} Bartoli, D., Faina, G., Marcugini, S., Pambianco, F. (2017). A construction of small complete caps in projective spaces. Journal of Geometry, 108, 215--246.
			
			\bibitem{evencap4} Bartoli, D., Giulietti, M., Marino, G., Polverino, O. (2017). Maximum scattered linear sets and complete caps in Galois spaces, Combinatorica, 38, 255--278.
			
			%\bibitem{Bartoli} Bartoli, D., Faina, G., Marcugini, S.,  Pambianco, F. (2017). A construction of small complete caps in projective spaces. Journal of Geometry, 108, 215--246.
			
			%\bibitem{affinecap} Bartoli, D., Faina, G., Marcugini, S., Pambianco, F. (2017). Complete caps in $\mathrm{AG}(N,q)$ with both $N$ and $q$ odd, J. Combin. Des. 25, 419--425.
			
			\bibitem{completecap} Cossidente, A., Csajb\'ok, B., Marino, G., Pavese, F. (2023). Small complete caps in $\mathrm{PG}(4n+1,q)$, Bull. Lond. Math. Soc., 55, 522--535.
			
			\bibitem{Chen} Chen, Y. G. (2018). On AP 3-covering sequences. Comptes Rendus. Mathématique, 356(2), 121--124.
			
			
			\bibitem{Cille}  Cilleruelo, J. (2012). Combinatorial problems in finite fields and Sidon sets. Combinatorica, 32(5), 497--511.
			
			\bibitem{handbook} Colbourne, C., Dinitz, J. (Eds.) (2007). Handbook of combinatorial designs. Boca Raton, FL: CRC press.
			
			\bibitem{CLP} Croot, E., Lev, V. F.,  Pach, P. P. (2017). Progression-free sets in $\mathbb{Z}_4^n$ are exponentially small. Annals of Mathematics, 185, 331--337.
			
			\bibitem{Pott} Czerwinski, I.,  Pott, A. (2023). Sidon sets, sum-free sets and linear codes. arXiv preprint arXiv:2304.07906.
			
			\bibitem{evencap3} Davydov, A.A., Giulietti, M., Marcugini, S., Pambianco, F. (2010), New inductive constructions of complete caps in $\mathrm{PG}(N,q)$, $q$ even, J. Combin. Des., 18, 177--201.
			
			%\bibitem{Dav} Davydov, A. A., Faina, G., Marcugini, S.,  Pambianco, F. (2009). On sizes of complete caps in projective spaces $\PG(n, q)$ and arcs in planes $\PG(2, q)$. Journal of Geometry, 94(1-2), 31-58.
			
			
			%\bibitem{Dav2} Davydov, A. A., Marcugini, S.,  Pambianco, F. (2004). Complete caps in projective spaces PG($n, q$). Journal of Geometry, 80, 23-30.
			\bibitem{DO} Davydov, A. A.,  \"Osterg\aa rd, P. R. (2001). Recursive constructions of complete caps. Journal of statistical planning and inference, 95(1-2), 167--173. 
			
			\bibitem{EB}
			Edel, Y.,  Bierbrauer, J. (2001). Large caps in small spaces. Designs, Codes and Cryptography, 23, 197--212.
			
			\bibitem{Bruyn} De Bruyn, J. V. D.,  Gijswijt, D. (2023). On the size of subsets of $\mathbb {F} _q^n $ avoiding solutions to linear systems with repeated columns.  Electronic Journal of Combinatorics, 30(4).
			
			\bibitem{Eber} Eberhard, S.,  Manners, F. (2023). The Apparent Structure of Dense Sidon Sets. The Electronic Journal of Combinatorics, P1-33.
			
			\bibitem{Edel}
			Edel, Y., Ferret, S., Landjev, I.,  Storme, L. (2002). The classification of the largest caps in $\AG(5, 3)$. Journal of Combinatorial Theory, Ser. A, 99(1), 95--110.
			
			\bibitem{EG}
			Ellenberg, J. S.,  Gijswijt, D. (2017). On large subsets of with no three-term arithmetic progression. Annals of Mathematics, 185, 339--343.
			
			\bibitem{EPPP}
			Elsholtz, C.,  Pach, P. P. (2020). Caps and progression-free sets in $\mathbb{Z}_m^n$. Designs, Codes and Cryptography, 88(10), 2133--2170.
			
			\bibitem{ET} Erd\H os, P.,  Turán, P. (1941). On a problem of Sidon in additive number theory, and on some related problems. J. London Math. Soc, 16(4), 212--215.
			
			\bibitem{Fang} Fang, J. H. (2021). A note on AP 3-covering sequences. Periodica Mathematica Hungarica, 83, 67--70.
			
			\bibitem{FGR}
			Frankl, P., Graham, R. L., Rödl, V. (1987). On subsets of abelian groups with no 3-term arithmetic progression. Journal of Combinatorial Theory, Ser. A, 45(1), 157--161.%\todo{erre szövegesen is hiv.}
			
			\bibitem{Fried} Fried, K.(1988). Rare bases for finite intervals of integers. Acta Sci. Math.(Szeged), 52(3–4), 303--305.
			
			\bibitem{evencap2} Giulietti, M. (2007). Small complete caps in $\mathrm{PG}(N,q)$, $q$ even, J. Combin. Des., 15, 420--436.
			
			\bibitem{Giuli}
			Giulietti, M. (2007). Small complete caps in Galois affine spaces. Journal of Algebraic Combinatorics, 25, 149--168.
			
			\bibitem{Nath}
			Güntürk, C. S.,  Nathanson, M. B. (2006). A new upper bound for finite additive bases. Acta Arithmetica, 124(3), 235--255.
			
			\bibitem{Habsi} Habsieger, L. (2014). On finite additive 2-bases. Transactions of the American Mathematical Society, 366(12), 6629--6646.
			
			\bibitem{Hirs} Hirschfeld, J. W.,  Storme, L. (2001, July). The packing problem in statistics, coding theory and finite projective spaces: update 2001. In Finite Geometries: Proceedings of the Fourth Isle of Thorns Conference (pp. 201--246). Boston, MA: Springer US.
			
			\bibitem{Hof} Hofmeister, G. (2001). Thin bases of order two. Journal of Number Theory, 86(1), 118--132.
			
			\bibitem{Tait} Huang, Y., Tait, M.,  Won, R. (2019). Sidon sets and 2-caps in $\F_3^n$. Involve, a Journal of Mathematics, 12(6), 995--1003.
			
			%\bibitem{Kara}
			%Karapetyan, K. I. (2022). Complete Caps in Affine Geometry AG(n,3). Mathematical Problems of Computer Science, 57, 56--64.
			
			\bibitem{SandorCs}   Kiss, S.Z.,   Sándor, Cs., Yang,  Q.H. (2018).  On generalized Stanley sequences. Acta Math. Hung., 154, 501--510.
			
			\bibitem{KovacsNagy} Kovács, B., Nagy, Z. L. (2025). Avoiding intersections of given size in finite affine spaces $\AG(n, 2)$. Journal of Combinatorial Theory, Ser. A, 209,
			105959.
			
			\bibitem{Meshulam}
			Meshulam, R. (1995). On subsets of finite abelian groups with no 3-term arithmetic progressions. Journal of Combinatorial Theory, Ser. A, 71(1), 168--172.
			
			\bibitem{Mimura} Mimura, M.,  Tokushige, N. (2021). Solving linear equations in a vector space over a finite field. Discrete Mathematics, 344(12), 112603.
			
			\bibitem{Mrose}  Mrose, A. (1979, April). Untere schranken für die reichweiten von extremalbasen fester ordnung. In Abhandlungen aus dem Mathematischen Seminar der Universität Hamburg (Vol. 48, pp. 118-124). Springer-Verlag.
			
			\bibitem{Nagy} Nagy, Z. L. (2018). Saturating sets in projective planes and hypergraph covers. Discrete Mathematics, 341(4), 1078--1083.
			
			\bibitem{GNagy} Nagy, G. P. (2022). Thin Sidon sets and the nonlinearity of vectorial Boolean functions. arXiv preprint arXiv:2212.05887.
			
			\bibitem{red} Redman, M., Rose, L.,  Walker, R. (2022). A Small Maximal Sidon Set in $\mathbb{Z}_2^n$. SIAM Journal on Discrete Mathematics, 36(3), 1861--1867.
			
			\bibitem{Ruzsa2} Ruzsa, I. Z. (1993). Solving a linear equation in a set of integers I. Acta Arithmetica, 65(3), 259--282.
			
			\bibitem{Ruzsa} Ruzsa, I. Z. (1998). A small maximal Sidon set. Analytic and Elementary Number Theory: A Tribute to Mathematical Legend Paul Erdős, 55--58.
			
			
			\bibitem{Oster} Östergard, P. R. (2000). Computer search for small complete caps. Journal of Geometry, 69(1-2), 172--179.
			
			\bibitem{evencap1} Pambianco, F. Storme, L. (1996). Small complete caps in spaces of even characteristci, Journal of Combinatorial Theory Ser. A.,  75, 70--84.
			
			%\bibitem{Platon} Platoni, I. (2014). Complete caps in $\AG(3, q)$ from elliptic curves. Journal of Algebra and Its Applications, 13(08), 1450050.
			
			\bibitem{Lisa} Sauermann, L. (2023). Finding solutions with distinct variables to systems of linear equations over $\F_p$. Mathematische Annalen, 386(1-2), 1-33.
			
			\bibitem{Ilja} Shkredov, I. D. (2006). Szemeredi's theorem and problems on arithmetic progressions. Russian Mathematical Surveys, 61(6), 1101.
			
			\bibitem{Smullyan} Smullyan, R. M. (1961). Theory of formal systems. Princeton University Press.
			
			\bibitem{Ty} Tyrrell, F. (2023). New lower bounds for cap sets,  Discrete Analysis, 20, 1-18.
			%\href{http://arxiv.org/abs/1512.07467}{http://arxiv.org/abs/1512.07467}
			
			%\bibtem{}
			
			
	\end{thebibliography}}
	
\end{document}
%\section{homokozó}


%\begin{itemize}
% \item valséges biz $[1,n]^k$-ra. Előtte kitalálni, nem lehet hogy ez az egész eleve ismert-e, monjuk $k=1$-re és onnan persze magasabb $\dim$-ra
%  \item megfogalmazás és kapcsolat: 3-\mathrm{AP}-k hipergráfjában domináló halmaz: elempárokkal domináljuk a hármasok 3. elemét. így lényegében domináló halmazt keresünk, amikor szaturálóról beszélünk. ( hypergraphs domination \cite{}) alter: $\beta$-ratioval lefedő ponthalmaz (a generált hármasok lefednek mindent.)
% \item leírás arról, hogy $\PG$-beli cap és $\AG$-beli cap közötti kapcsolatot hogyan lehet megfogni. ha vki $\PG$-beli cap, és ráadásul elkerül egy hipersíkot, akkor ő $\AG$-beli cap is. Davydov konstrukciói közül így van ami jól működik (lasd. egyik Giulietti cikk egyik megjegyzese).

%   \item Van e $cq$-s complete $\beta$-ratio-free $\F_q^2$-ben ha $\beta$ ès $-\beta-1$ n\'egyzet? 
%  \item a véletlen módszer megokosítása: az iterálva kiválasztott $H_i$ halmaz által adott  fedett pontok, amikből nem akarunk választani, valójában a nem-fedettekből ugyanannyit fednének kb, mint a nem fedett pontok generátuma, várható értékben.
%   \item 1-r\H ol meg $U$-avoidingr\'ol \'att\'erni $\beta$-ra.
%  \item Theorem 2.3. $\beta$-ra $\F_q$ (milyen felt\'etel lesz $q$-ra $\beta$ fv.-\'eben)
%   \item Result 3.8, Corollary 3.9 $\beta$-ra? $\F_q \times \F_q$ helyett masutt is ad-e ilyen szepet ez a konstrukcio.
%  \item $\AG(3,q)$-ban k\'et hipersikbol elhagyni pontokat? Pl. egy hipersik es egy masik hipersikot szaturalo halmaz jo vagy jova teheto de ez nagy.
%    \item Sidon -> 3AP free Mimura, M.,\& Tokushige, N. (2021). Solving linear equations in a vector space over a finite field. Discrete Mathematics, 344(12), 112603.
%  \item \url{https://eudml.org/doc/206579}
%   \item \url{https://link.springer.com/article/10.1007/s00208-022-02391-y}
%  \item On complete Sidons in $\F_2^t$ by Pott et al \todo{very recent :-/ }:
%   \url{https://arxiv.org/pdf/2304.07906.pdf}
%  Is there any chance that the parabola is sometimes complete Sidon? 
%   \item $[1,n]\times[1,n]$-ben Costas Array!!
%   \item small complete Sidon over $\F_2$: \url{https://arxiv.org/pdf/2109.00292.pdf}
%  \item Ruzsa: a small maximal sidon \url{https://link.springer.com/article/10.1023/A:1009757824153}
%   \item Nagy Gabortol ez jo, complete kis Sidonok: \url{https://arxiv.org/pdf/2212.05887.pdf}
%  \item For product of caps: Edel and the references there: Extensions of Generalized Product Caps
%   \item tavolrol nezve a 4.fejezetben mintha vmi olyasmit csinalna mint mi a parabolaval
%\url{https://www.sciencedirect.com/science/article/pii/S0022314X17303578}
%  \item Ebben nagyon sok "nagy" Sidon ossze van szedve, es ezeknel van r\'a esely, hogy esetleg minimalis meretu complete $3$-$\mathrm{AP}$-free-k, kulonbozo csoportokban. Jo lenne ezeket atnezni.
%\end{itemize}




